\beginnumbering
\eledsection*{Adhyāya 27}
\pstart \ledsidenote{\passref{6.27.1}}\skt{अथातो }\edtext{\edtext{\skt{नव\-ग्र\-ह\-वि\-ज्ञा\-नी\-य\-म\-ध्यायं}}{\lemma{\skt{नव\appbrk{}ग्र\appbrk{}ह\appbrk{}वि\appbrk{}ज्ञा\appbrk{}नी\appbrk{}यम्}}\Afootnote{K; \skt{नव\appbrk{}ग्र\appbrk{}हा\appbrk{}कृ\appbrk{}ति\appbrk{}वि\appbrk{}ज्ञा\appbrk{}नी\appbrk{}यम्}\ A; \skt{नवग्रहविज्ञा}\ H}}}{\lemma{\skt{अध्यायं}}\Afootnote{A K; \skt{यम्}\ H}}\skt{ }\edtext{\skt{व्याख्यास्यामः}}{\lemma{\skt{व्याख्यास्यामः}}\Afootnote{A H; \emph{om.}\ K}}\pend

\pstart\versepar \ledsidenote{\passref{6.27.3}}\edtext{\skt{बाल\-ग्र\-हा\-णा\-म्वि\-ज्ञानं}}{\lemma{\skt{बालग्रहाणाम्}}\Afootnote{A H; \skt{वालग्रहानाम्}\ K}}\skt{ }\edtext{\skt{साधणां}}{\lemma{\skt{साधणां}}\Afootnote{K; \skt{साधनं}\ A H}}\skt{ }\edtext{\skt{चाप्यनन्तरं}}{\lemma{\skt{चाप्यनन्तरं}}\Afootnote{K H; \skt{चाप्यनन्तर}\ A}}\skt{~। }\pend
\pstart\versepar \skt{}\edtext{\skt{उत्पत्तिं}}{\lemma{\skt{उत्पत्तिं}}\Afootnote{K H; \skt{म्~। उत्पत्तिं}\ A}}\skt{ }\edtext{\skt{कारणां}}{\lemma{\skt{कारणां}}\Afootnote{K; \skt{कारणं}\ A; \skt{कारणैञ्}\ H}}\skt{ चैव }\edtext{\skt{सुश्रुतैकमना}}{\lemma{\skt{सुश्रुतैकमना}}\Afootnote{\skt{सुश्रुतैकमनाः}\ A; \skt{सश्रुतैकमना}\ K H}}\skt{ शृणु~।}\pend

\pstart\versepar \ledsidenote{\passref{6.27.4}}\skt{स्कन्दग्रहस्तु प्रथमः स्कन्दापस्मार एव च~। }\pend
\pstart\versepar \skt{शकुनी रेवती चैव पूतना चान्धपूतना~।}\pend

\pstart\versepar \ledsidenote{\passref{6.27.5}}\skt{पूतना शीतनामा च तथैव मुखमण्डिका~। }\pend
\pstart\versepar \skt{नवमो नैगमेषश्च यः }\edtext{\skt{पितृ\-ग्र\-ह\-सं\-ज्ञितः}}{\lemma{\skt{पितृ\appbrk{}ग्र\appbrk{}ह\appbrk{}सं\appbrk{}ज्ञितः}}\Afootnote{A; \skt{पितृ\appbrk{}ग्र\appbrk{}ह\appbrk{}सं\appbrk{}ज्ञिनः}\ K; \skt{पितृ\appbrk{}ग्र\appbrk{}ह\appbrk{}सं\appbrk{}ज्ञिताः}\ H}}\skt{~॥}\pend

\pstart\versepar \ledsidenote{\passref{6.27.6}}\edtext{\skt{धात्री\-मा\-त्रो\-श्चौ\-प\-चा\-रा\-प\-र\-द्वाञ्}}{\lemma{\skt{धात्री\appbrk{}मा\appbrk{}त्रो\appbrk{}श्चौ\appbrk{}प\appbrk{}चा\appbrk{}रा\appbrk{}प\appbrk{}र\appbrk{}द्वाञ्}}\Afootnote{K; \skt{धात्रीमात्रोः प्राक्प्र\appbrk{}दि\appbrk{}ष्टा\appbrk{}प\appbrk{}चारा}\ A; \skt{धात्री\appbrk{}मा\appbrk{}त्रो\appbrk{}श्चो\appbrk{}प\appbrk{}चा\appbrk{}रा\appbrk{}प\appbrk{}र\appbrk{}द्वाञ्}\ H}}\skt{ }\pend
\pstart\versepar \skt{}\edtext{\skt{छौचभ्रष्टं}}{\lemma{\skt{छौचभ्रष्टं}}\Afootnote{K H; \skt{च्छौचभ्रष्टान्}\ A}}\skt{ मङ्ग\-ला\-चा\-र\-हीनां~। }\pend
\pstart\versepar \skt{}\edtext{\skt{त्रस्तां}}{\lemma{\skt{त्रस्तां}}\Afootnote{K H; \skt{त्रस्ता}\ A}}\skt{ }\edtext{\skt{भीतां}}{\lemma{\skt{भीतां}}\Afootnote{K H; \skt{न्हृष्टां}\ A}}\skt{ }\edtext{\skt{त्रिताञ्स्तर्जितान्वा}}{\lemma{\skt{त्रिताञ्स्तर्जितान्वा}}\Afootnote{K H; \skt{तर्जि\appbrk{}ता\appbrk{}न्ता\appbrk{}डि\appbrk{}तान्वा}\ A}}\skt{ }\pend
\pstart\versepar \skt{}\edtext{\skt{पूजा हि स्युरेते}}{\lemma{\skt{पूजा हि स्युरेते}}\Afootnote{\skt{पूजा\appbrk{}हे\appbrk{}तो\appbrk{}र्हिं\appbrk{}स्यु\appbrk{}रेते}\ A; \skt{पूजाहेतोर्हि स्युरेते}\ K H}}\skt{ }\edtext{\skt{कुमारां}}{\lemma{\skt{कुमारां}}\Afootnote{K; \skt{कुमारान्}\ A H}}\skt{~।}\pend

\pstart\versepar \ledsidenote{\passref{6.27.7}}\skt{ऐश्वर्यस्थास्ते न }\edtext{\edtext{\skt{शक्यम्विशन्तो}}{\lemma{\skt{शक्यम्}}\Afootnote{K H; \emph{om.}\ A}}}{\lemma{\skt{विशन्तो}}\Afootnote{K; \skt{शक्या विशन्तो}\ A; \skt{विशक्यं विशन्तो}\ H}}\skt{ }\edtext{\skt{देअं}}{\lemma{\skt{देअं}}\Afootnote{K; \skt{देहं}\ A; \skt{वेदं}\ H}}\skt{ द्रष्टुं }\edtext{\skt{मानु\-षै\-र्वि\-श्व\-रूपा}}{\lemma{\skt{मानु\appbrk{}षै\appbrk{}र्वि\appbrk{}श्व\appbrk{}रूपा}}\Afootnote{K; \skt{मानु\appbrk{}षै\appbrk{}र्वि\appbrk{}श्व\appbrk{}रूपाः}\ A H}}\skt{~। }\pend
\pstart\versepar \skt{आप्तं वाक्यं }\edtext{\skt{सम्य\-गी\-क्ष्या\-भि\-धास्ये}}{\lemma{\skt{सम्य\appbrk{}गी\appbrk{}क्ष्या\appbrk{}भि\appbrk{}धास्ये}}\Afootnote{K H; \skt{तत्स\appbrk{}मी\appbrk{}क्ष्या\appbrk{}भि\appbrk{}धास्ये}\ A}}\skt{ लिङ्गान्येषां यानि }\edtext{\skt{देहेषु}}{\lemma{\skt{देहेषु}}\Afootnote{K H; \skt{देहे भ}\ A}}\skt{ }\edtext{\skt{तानि}}{\lemma{\skt{तानि}}\Afootnote{K H; \skt{वन्ति}\ A}}\skt{~।}\pend

\pstart\versepar \ledsidenote{\passref{6.27.8}}\edtext{\skt{स्तब्धाक्षः}}{\lemma{\skt{स्तब्धाक्षः}}\Afootnote{K H; \skt{शूनाक्षः}\ A}}\skt{ }\edtext{\skt{क्षत\-व\-स\-ग\-न्धिकः}}{\lemma{\skt{क्षत\appbrk{}व\appbrk{}स\appbrk{}ग\appbrk{}न्धिकः}}\Afootnote{K H; \skt{क्षत\appbrk{}ज\appbrk{}स\appbrk{}ग\appbrk{}न्धिकः}\ A}}\skt{ }\edtext{\skt{स्तनद्विड्वक्रास्यो}}{\lemma{\skt{स्तनद्विड्वक्रास्यो}}\Afootnote{A K; \skt{स्तनद्विड् {\symfont ¦}स्यो}\ H}}\skt{ }\edtext{\skt{हत\-च\-र\-णे\-क\-प\-त्र\-नेत्रः}}{\lemma{\skt{हत\appbrk{}च\appbrk{}र\appbrk{}णे\appbrk{}क\appbrk{}प\appbrk{}त्र\appbrk{}नेत्रः}}\Afootnote{K H; \skt{हत\appbrk{}च\appbrk{}लि\appbrk{}तै\appbrk{}क\appbrk{}प\appbrk{}क्ष्म\appbrk{}नेत्रः}\ A}}\skt{~। }\pend
\pstart\versepar \skt{}\edtext{\skt{उद्विग्नः}}{\lemma{\skt{उद्विग्नः}}\Afootnote{K H; \skt{उद्विग्रः}\ A}}\skt{ }\edtext{\skt{सलु\-लि\-त\-च\-क्षु\-र\-ल्प\-रोगी}}{\lemma{\skt{सलु\appbrk{}लि\appbrk{}त\appbrk{}च\appbrk{}क्षु\appbrk{}र\appbrk{}ल्प\appbrk{}रोगी}}\Afootnote{K H; \skt{सुलु\appbrk{}लि\appbrk{}त\appbrk{}च\appbrk{}क्षु\appbrk{}र\appbrk{}ल्प\appbrk{}रोदी}\ A}}\skt{ स्कन्दार्त्तो भवति च }\edtext{\skt{गाढमुष्टिवर्धा}}{\lemma{\skt{गाढमुष्टिवर्धा}}\Afootnote{K; \skt{गाढमुष्टिवर्चाः}\ A; \skt{गाढमुष्टिवन्धा}\ H}}\skt{~।}\pend

\pstart\versepar \ledsidenote{\passref{6.27.9}}\edtext{\skt{निःसंज्ञ}}{\lemma{\skt{निःसंज्ञ}}\Afootnote{K; \skt{निःसंज्ञो}\ A H}}\skt{ भवति }\edtext{\skt{पुनर्लभेत्च}}{\lemma{\skt{पुनर्लभेत्च}}\Afootnote{K H; \skt{पुनर्भवेत्स}\ A}}\skt{ }\edtext{\skt{सञ्ज्ञं}}{\lemma{\skt{सञ्ज्ञं}}\Afootnote{K H; \skt{संज्ञः}\ A}}\skt{ संरब्धः करचरणैश्च नृत्यतीव~। }\pend
\pstart\versepar \skt{}\edtext{\skt{विण्मूत्रसृजति}}{\lemma{\skt{विण्मूत्रसृजति}}\Afootnote{K; \skt{विण्मूत्रे सृजति}\ A; \skt{विण्मूत्रैः सृजति}\ H}}\skt{ विनद्य }\edtext{\skt{जृम्भमाणं}}{\lemma{\skt{जृम्भमाणं}}\Afootnote{K H; \skt{जृम्भमाणः}\ A}}\skt{ }\edtext{\skt{फेनाञ्श्च}}{\lemma{\skt{फेनाञ्श्च}}\Afootnote{K H; \skt{फेनं च}\ A}}\skt{ }\edtext{\skt{सृजति च}}{\lemma{\skt{सृजति च}}\Afootnote{K H; \skt{प्रसृजति}\ A}}\skt{ तत्सखाभिपन्नः~॥}\pend

\pstart\versepar \ledsidenote{\passref{6.27.10}}\skt{स्रस्तांगो भयचकितो }\edtext{\skt{विहंगगन्धा सा}}{\lemma{\skt{विहंगगन्धा सा}}\Afootnote{K H; \skt{विहङ्गगन्धिः सं}\ A}}\skt{ }\edtext{\skt{स्राव\-व्र\-ण\-प\-रि\-पी\-डितः}}{\lemma{\skt{स्राव\appbrk{}व्र\appbrk{}ण\appbrk{}प\appbrk{}रि\appbrk{}पी\appbrk{}डितः}}\Afootnote{K; \skt{स्रावि\appbrk{}व्र\appbrk{}ण\appbrk{}प\appbrk{}रि\appbrk{}पी\appbrk{}डितः}\ A; \skt{स्राव\appbrk{}व्र\appbrk{}ण\appbrk{}प\appbrk{}रि\appbrk{}प\appbrk{}रि\appbrk{}पी\appbrk{}डितः}\ H}}\skt{ समन्तात्~। }\pend
\pstart\versepar \skt{स्फोटैश्च प्रचिततनुः }\edtext{\skt{सदा\-ह\-पा\-कै\-र्निः\-संज्ञो}}{\lemma{\skt{सदा\appbrk{}ह\appbrk{}पा\appbrk{}कै\appbrk{}र्निः\appbrk{}संज्ञो}}\Afootnote{K H; \skt{सदा\appbrk{}ह\appbrk{}पा\appbrk{}कै\appbrk{}र्वि\appbrk{}ज्ञेयो}\ A}}\skt{ भवति शिशुः क्षतः शकुन्या~॥}\pend

\pstart\versepar \ledsidenote{\passref{6.27.11}}\skt{रक्तास्यो }\edtext{\skt{हरि\-त\-म\-लो\-भि\-पा\-ण्डु\-देहः}}{\lemma{\skt{हरि\appbrk{}त\appbrk{}म\appbrk{}लो\appbrk{}भि\appbrk{}पा\appbrk{}ण्डु\appbrk{}देहः}}\Afootnote{K H; \skt{हरि\appbrk{}त\appbrk{}म\appbrk{}लोऽ\appbrk{}ति\appbrk{}पा\appbrk{}ण्डु\appbrk{}देहः}\ A}}\skt{ श्यावो वा }\edtext{\skt{ज्वर\-मु\-ख\-पा\-क\-वे\-द\-नार्त्तः}}{\lemma{\skt{ज्वर\appbrk{}मु\appbrk{}ख\appbrk{}पा\appbrk{}क\appbrk{}वे\appbrk{}द\appbrk{}नार्त्तः}}\Afootnote{K H; \skt{ज्वर\appbrk{}मु\appbrk{}ख\appbrk{}पा\appbrk{}क\appbrk{}वे\appbrk{}द\appbrk{}नातः}\ A}}\skt{~। }\pend
\pstart\versepar \skt{}\edtext{\skt{मृद्रीते}}{\lemma{\skt{मृद्रीते}}\Afootnote{K H; \skt{रेवत्या वि}\ A}}\skt{ }\edtext{\skt{व्याथिततनुश्च}}{\lemma{\skt{व्याथिततनुश्च}}\Afootnote{K H; \skt{यथिततनुश्च}\ A}}\skt{ }\edtext{\skt{कर्णनासां}}{\lemma{\skt{कर्णनासां}}\Afootnote{K H; \skt{कर्णनासं}\ A}}\skt{ }\edtext{\skt{रेवत्या}}{\lemma{\skt{रेवत्या}}\Afootnote{K H; \skt{मृद्गाति}\ A}}\skt{ }\edtext{\skt{भृश\-म\-ति\-पी\-डितः}}{\lemma{\skt{भृश\appbrk{}म\appbrk{}ति\appbrk{}पी\appbrk{}डितः}}\Afootnote{K H; \skt{ध्रुव\appbrk{}म\appbrk{}भि\appbrk{}पी\appbrk{}डितः}\ A}}\skt{ कुमारः~॥}\pend

\pstart\versepar \ledsidenote{\passref{6.27.12}}\skt{स्रस्ताङ्गः स्वपिति सुखं }\edtext{\skt{दिनरात्रौ}}{\lemma{\skt{दिनरात्रौ}}\Afootnote{K; \skt{दिवा न रात्रौ}\ A H}}\skt{ }\edtext{\skt{विन्नां}}{\lemma{\skt{विन्नां}}\Afootnote{K; \skt{विड्भिन्नं}\ A; \skt{विभिन्नां}\ H}}\skt{ सृजति च काकतुल्यगन्धिः~। }\pend
\pstart\versepar \skt{छर्द्यार्तो हृषि\-त\-त\-नू\-रुहः कुमा\-र\-स्तृ\-ष्णा\-लु\-र्भ\-वति च पूतनागृहीतः~॥}\pend

\pstart\versepar \ledsidenote{\passref{6.27.13}}\skt{यो द्वेष्टि }\edtext{\skt{स्तन\-म\-ति\-सा\-र\-का\-स\-हि\-क्का\-च्छ\-र्द्दिर्वा}}{\lemma{\skt{स्तन\appbrk{}म\appbrk{}ति\appbrk{}सा\appbrk{}र\appbrk{}का\appbrk{}स\appbrk{}हि\appbrk{}क्का\appbrk{}च्छ\appbrk{}र्द्दिर्वा}}\Afootnote{K H; \skt{स्तन\appbrk{}म\appbrk{}ति\appbrk{}सा\appbrk{}र\appbrk{}का\appbrk{}स\appbrk{}हि\appbrk{}क्का\appbrk{}च्छ\appbrk{}र्दीभि}\ A}}\skt{ }\edtext{\skt{ज्वर\-स\-हि\-ता\-भि\-र\-र्द्य\-मानः}}{\lemma{\skt{ज्वर\appbrk{}स\appbrk{}हि\appbrk{}ता\appbrk{}भि\appbrk{}र\appbrk{}र्द्य\appbrk{}मानः}}\Afootnote{K H; \skt{र्ज्वर\appbrk{}स\appbrk{}हि\appbrk{}ता\appbrk{}भि\appbrk{}र\appbrk{}र्द्य\appbrk{}मानः}\ A}}\skt{~। }\pend
\pstart\versepar \skt{}\edtext{\skt{दुर्वर्णाः}}{\lemma{\skt{दुर्वर्णाः}}\Afootnote{K; \skt{दुर्वर्णः}\ A H}}\skt{ }\edtext{\skt{सततमथापि}}{\lemma{\skt{सततमथापि}}\Afootnote{K H; \skt{सततमधःश}\ A}}\skt{ }\edtext{\skt{चाम्लगन्धीतं}}{\lemma{\skt{चाम्लगन्धीतं}}\Afootnote{K H; \skt{योऽम्लगन्धिस्तं}\ A}}\skt{ }\edtext{\skt{यादु\-प\-ह\-त\-म\-न्ध\-पू\-त\-नायाः}}{\lemma{\skt{यादु\appbrk{}प\appbrk{}ह\appbrk{}त\appbrk{}म\appbrk{}न्ध\appbrk{}पू\appbrk{}त\appbrk{}नायाः}}\Afootnote{K; \skt{ब्रूयुर्भिषज इहा\appbrk{}न्ध\appbrk{}पू\appbrk{}त\appbrk{}ना\appbrk{}र्तम्}\ A; \skt{व्व्रूया\appbrk{}दु\appbrk{}प\appbrk{}ह\appbrk{}त\appbrk{}म\appbrk{}न्ध\appbrk{}पू\appbrk{}त\appbrk{}नायाः}\ H}}\skt{~।}\pend

\pstart\versepar \ledsidenote{\passref{6.27.14}}\edtext{\skt{आक्रन्दत्यति}}{\lemma{\skt{आक्रन्दत्यति}}\Afootnote{K H; \skt{उद्विग्रो भृ}\ A}}\skt{ }\edtext{\skt{चकितस्तु}}{\lemma{\skt{चकितस्तु}}\Afootnote{K H; \skt{शमति}\ A}}\skt{ }\edtext{\skt{वेपमानः}}{\lemma{\skt{वेपमानः}}\Afootnote{K H; \skt{वेपते प्ररुद्या}\ A}}\skt{ }\edtext{\skt{संलीनः}}{\lemma{\skt{संलीनः}}\Afootnote{K H; \skt{त्संलीनः स्}\ A}}\skt{ स्वपिति च यस्य चान्त्र }\edtext{\skt{कूजाः}}{\lemma{\skt{कूजाः}}\Afootnote{K; \skt{कूजः}\ A; \skt{कूजा}\ H}}\skt{~। }\pend
\pstart\versepar \skt{}\edtext{\skt{स्रस्ताङ्गो}}{\lemma{\skt{स्रस्ताङ्गो}}\Afootnote{K H; \skt{विस्राङ्गो}\ A}}\skt{ }\edtext{\skt{भुश\-म\-ति\-सा\-र्यते यः}}{\lemma{\skt{भुश\appbrk{}म\appbrk{}ति\appbrk{}सा\appbrk{}र्यते यः}}\Afootnote{K; \skt{भृश\appbrk{}म\appbrk{}ति\appbrk{}सा\appbrk{}र्यते च}\ A H}}\skt{ }\edtext{\skt{स्यात्तं}}{\lemma{\skt{स्यात्तं}}\Afootnote{K; \skt{यस्तं}\ A; \skt{यो वै}\ H}}\skt{ }\edtext{\skt{कुर्या\-द्भि\-ष\-ग\-थ\-शी\-त\-पू\-त\-नार्तः}}{\lemma{\skt{कुर्या\appbrk{}द्भि\appbrk{}ष\appbrk{}ग\appbrk{}थ\appbrk{}शी\appbrk{}त\appbrk{}पू\appbrk{}त\appbrk{}नार्तः}}\Afootnote{K; \skt{जानी\appbrk{}या\appbrk{}द्भि\appbrk{}ष\appbrk{}गिह शीत\appbrk{}पू\appbrk{}त\appbrk{}ना\appbrk{}र्तम्}\ A; \skt{तं व्रूया भिषगथ शीतपूतनार्तः}\ H}}\skt{~।}\pend

\pstart\versepar \ledsidenote{\passref{6.27.15}}\skt{म्लानाङ्गः }\edtext{\skt{सुरु\-चि\-र\-पा\-णि\-पा\-द\-वक्रो}}{\lemma{\skt{सुरु\appbrk{}चि\appbrk{}र\appbrk{}पा\appbrk{}णि\appbrk{}पा\appbrk{}द\appbrk{}वक्रो}}\Afootnote{K H; \skt{सुरु\appbrk{}चि\appbrk{}र\appbrk{}पा\appbrk{}णि\appbrk{}पा\appbrk{}द\appbrk{}वक्त्रो}\ A}}\skt{ }\edtext{\skt{वह्वाशी}}{\lemma{\skt{वह्वाशी}}\Afootnote{K H; \skt{बह्वाशी}\ A}}\skt{ कलु\-ष\-सि\-रा\-वृ\-तो\-दरो यः~। }\pend
\pstart\versepar \skt{}\edtext{\skt{संसृष्टः}}{\lemma{\skt{संसृष्टः}}\Afootnote{K H; \emph{om.}\ A}}\skt{ }\edtext{\skt{शिशु\-र\-थ\-व\-क्रु\-म\-ण्डि\-कायाः}}{\lemma{\skt{शिशु\appbrk{}र\appbrk{}थ\appbrk{}व\appbrk{}क्रु\appbrk{}म\appbrk{}ण्डि\appbrk{}कायाः}}\Afootnote{K; \emph{om.}\ A; \skt{शिशुरथ वक्रुमण्डिका या}\ H}}\skt{ सोद्वेगो भवति च }\edtext{\skt{वस्तुमूत्रगन्धः}}{\lemma{\skt{वस्तुमूत्रगन्धः}}\Afootnote{K H; \skt{मूत्रतुल्यगन्धिः स ज्ञेयः शिशुरिह वक्त्रमण्डिकार्तः}\ A}}\skt{~॥}\pend

\pstart\versepar \ledsidenote{\passref{6.27.16}}\edtext{\skt{यः फेनाम्वमति}}{\lemma{\skt{यः फेनाम्}}\Afootnote{K; \skt{यः फेनं}\ A H}}\skt{ }\edtext{\skt{विनाम्यते च}}{\lemma{\skt{विनाम्यते च}}\Afootnote{K H; \skt{विनम्यते च}\ A}}\skt{ मध्ये }\edtext{\skt{सोद्वेगो}}{\lemma{\skt{सोद्वेगो}}\Afootnote{K H; \skt{सोद्वेगं}\ A}}\skt{ }\edtext{\skt{विहसति}}{\lemma{\skt{विहसति}}\Afootnote{K H; \skt{विलपति}\ A}}\skt{ }\edtext{\skt{चोर्ध्वमीक्षमानः}}{\lemma{\skt{चोर्ध्वमीक्षमानः}}\Afootnote{K H; \skt{चोर्ध्वमीक्षमाणः}\ A}}\skt{~। }\pend
\pstart\versepar \skt{}\edtext{\skt{ज्वर्येच्च}}{\lemma{\skt{ज्वर्येच्च}}\Afootnote{K H; \skt{ज्वर्येत}\ A}}\skt{ प्रततमथो }\edtext{\skt{वसाश्रगन्धी}}{\lemma{\skt{वसाश्रगन्धी}}\Afootnote{K H; \skt{वसासगन्धि}\ A}}\skt{ }\edtext{\skt{सज्ञेयो}}{\lemma{\skt{सज्ञेयो}}\Afootnote{K H; \skt{र्निःसंज्ञो}\ A}}\skt{ भवति हि }\edtext{\skt{नैगमेषदष्टः}}{\lemma{\skt{नैगमेषदष्टः}}\Afootnote{K H; \skt{नैगमेषजुष्टः}\ A}}\skt{~॥}\pend

\pstart\versepar \ledsidenote{\passref{6.27.17}}\edtext{\skt{धात्री\-स्त\-न\-द्वि\-द्वि\-ष्टव}}{\lemma{\skt{धात्री\appbrk{}स्त\appbrk{}न\appbrk{}द्वि\appbrk{}द्वि\appbrk{}ष्टव}}\Afootnote{K; \skt{प्रस्तब्धो यः स्तनद्वे}\ A; \skt{धात्री\appbrk{}स्त\appbrk{}न\appbrk{}द्वि\appbrk{}द्वि\appbrk{}ष्टव्व}\ H}}\skt{ }\edtext{\skt{आविष्टो}}{\lemma{\skt{आविष्टो}}\Afootnote{K H; \skt{षी}\ A}}\skt{ मुह्यते मुहुः~। }\pend
\pstart\versepar \skt{तं }\edtext{\skt{बालं न}}{\lemma{\skt{बालं न}}\Afootnote{K H; \skt{बालम}\ A}}\skt{ }\edtext{\skt{चिराद्वन्ति}}{\lemma{\skt{चिराद्वन्ति}}\Afootnote{K H; \skt{चिराद्धन्ति}\ A}}\skt{ ग्रहः }\edtext{\skt{सम्पूर्णलक्षणं}}{\lemma{\skt{सम्पूर्णलक्षणं}}\Afootnote{K H; \skt{संपूर्णलक्षणः}\ A}}\skt{~॥}\edtext{}{\lemma{\skt{मुहुः~। तं}}\Afootnote{K H; \skt{चाविशन्मुहुः~। तं}\ A}}\pend

\pstart\versepar \ledsidenote{\passref{6.27.18}}\skt{विपरीतमतः साध्यं चिकि\-त्से\-द\-चि\-रा\-र्दितं~। }\pend
\pstart\versepar \skt{}\edtext{\skt{गृहं}}{\lemma{\skt{गृहं}}\Afootnote{K H; \skt{गृहे}\ A}}\skt{ पुरा\-ण\-ह\-वि\-षा\-भ्यज्य बालं शुचौ शुचिः~।}\pend

\pstart\versepar \ledsidenote{\passref{6.27.19}}\edtext{\skt{सर्षपा}}{\lemma{\skt{सर्षपा}}\Afootnote{K H; \skt{सर्षपान्}\ A}}\skt{ }\edtext{\skt{प्रकिरेत्तत्र}}{\lemma{\skt{प्रकिरेत्तत्र}}\Afootnote{K H; \skt{प्रकिरेत्तेषां}\ A}}\skt{ }\edtext{\skt{तैलदीपञ्च}}{\lemma{\skt{तैलदीपञ्च}}\Afootnote{K H; \skt{तैलैर्दीपं च}\ A}}\skt{ कारयेत्~। }\pend
\pstart\versepar \skt{सदा }\edtext{\skt{सन्निहितञ्चात्र}}{\lemma{\skt{चात्र}}\Afootnote{K; \skt{चापि}\ A; \skt{तु}\ H}}\skt{ }\edtext{\skt{जुहु\-या\-द्ध्व्य\-वा\-हनं}}{\lemma{\skt{जुहु\appbrk{}या\appbrk{}द्ध्व्य\appbrk{}वा\appbrk{}हनं}}\Afootnote{K H; \skt{जुहु\appbrk{}या\appbrk{}द्ध\appbrk{}व्य\appbrk{}वा\appbrk{}ह\appbrk{}नम्}\ A}}\skt{~।}\pend

\pstart\versepar \ledsidenote{\passref{6.27.20}}\edtext{\edtext{\skt{सर्व\-र\-नौ\-ष\-धी\-न्वी\-जा\-ग\-न्ध\-पु\-ष्पै\-र\-ल\-ङ्कृतां}}{\lemma{\skt{सर्व\appbrk{}र\appbrk{}नौ\appbrk{}ष\appbrk{}धीन्}}\Afootnote{K; \skt{सर्वगन्धौषधी}\ A; \skt{सर्व्व\appbrk{}व\appbrk{}नौ\appbrk{}ष\appbrk{}धीन्}\ H}}}{\lemma{\skt{वीजा\appbrk{}ग\appbrk{}न्ध\appbrk{}पु\appbrk{}ष्पै\appbrk{}र\appbrk{}ल\appbrk{}ङ्कृतां}}\Afootnote{K; \skt{बीजै\appbrk{}र्ग\appbrk{}न्ध\appbrk{}मा\appbrk{}ल्यै\appbrk{}र\appbrk{}ल\appbrk{}ङ्कृत}\ A; \skt{वीजा\appbrk{}न्ग\appbrk{}न्ध\appbrk{}पु\appbrk{}ष्पै\appbrk{}र\appbrk{}ल\appbrk{}ङ्कता}\ H}}\skt{~। }\pend
\pstart\versepar \skt{}\edtext{\skt{अग्नये}}{\lemma{\skt{अग्नये}}\Afootnote{K; \skt{म्~। अग्न्ये}\ A; \skt{न्~। अग्नये}\ H}}\skt{ कृत्तिकाभ्यश्च }\edtext{\skt{स्वाहा}}{\lemma{\skt{स्वाहा}}\Afootnote{A K; \emph{om.}\ H}}\skt{ स्वाहेति }\edtext{\skt{संततं}}{\lemma{\skt{संततं}}\Afootnote{A K; \skt{सततं}\ H}}\skt{~।}\pend

\pstart \ledsidenote{\passref{6.27.21}}\skt{नमस्स्कन्दाय देवाय ग्रहाधिपतये नमः~। शिरसा }\edtext{\skt{त्वाभिवन्दे हं}}{\lemma{\skt{त्वाभिवन्दे हं}}\Afootnote{A K; \skt{त्वभिवन्देहं}\ H}}\skt{ }\edtext{\skt{प्रतिगृण्हीष्व मे}}{\lemma{\skt{प्रतिगृण्हीष्व मे}}\Afootnote{K H; \skt{प्रतिगृह्णीष्व मे}\ A}}\skt{ बलिं~। }\edtext{\skt{नीरु\-जा\-नि\-र्व्वि\-का\-रश्च}}{\lemma{\skt{नीरु\appbrk{}जा\appbrk{}नि\appbrk{}र्व्वि\appbrk{}का\appbrk{}रश्च}}\Afootnote{K H; \skt{नीरुजो निर्विकारश्च}\ A}}\skt{ }\edtext{\skt{शिशुर्वाश्च}}{\lemma{\skt{शिशुर्वाश्च}}\Afootnote{K H; \skt{शिशुर्मे}\ A}}\skt{ }\edtext{\skt{जयतामिति}}{\lemma{\skt{जयतामिति}}\Afootnote{K; \skt{जायतां द्रुतम्}\ A; \skt{जायतामिति~॥ {\symfont ❈}}\ H}}\skt{~॥}\pend

\pstart \ledsidenote{\passref{6.27.trailer}}\edtext{\skt{कुमारे}}{\lemma{\skt{कुमारे}}\Afootnote{K; \skt{इति सुश्रु\appbrk{}त\appbrk{}सं\appbrk{}हि\appbrk{}ता\appbrk{}या\appbrk{}मु\appbrk{}त्त\appbrk{}र\appbrk{}त\appbrk{}न्त्रा\appbrk{}न्त\appbrk{}र्गते कुमारतन्त्रे नव\appbrk{}ग्र\appbrk{}हा\appbrk{}कृ\appbrk{}ति\appbrk{}वि\appbrk{}ज्ञा\appbrk{}नीयो नाम (प्रथमोऽध्यायः, आदितः) सप्तविंशोऽध्यायः}\ A; \skt{कुमारतन्त्रे नवग्रह{\symfont ¦}विज्ञानीयं प्रथमो ध्यायः}\ H}}\skt{~॥ १~॥ ०~॥}\pend

\endnumbering
\beginnumbering
\eledsection*{Adhyāya 28}
\pstart \ledsidenote{\passref{6.28.1}}\edtext{\skt{अथातः}}{\lemma{\skt{अथातः}}\Afootnote{A K; \skt{अथातो}\ H}}\skt{ }\edtext{\skt{स्कन्द\-ग्र\-ह\-प्र\-ति\-षे\-ध\-म\-ध्यायं}}{\lemma{\skt{स्कन्द\appbrk{}ग्र\appbrk{}ह\appbrk{}प्र\appbrk{}ति\appbrk{}षे\appbrk{}ध\appbrk{}म\appbrk{}ध्यायं}}\Afootnote{K; \skt{स्कन्ध\appbrk{}ग्र\appbrk{}ह\appbrk{}प्र\appbrk{}ति\appbrk{}षेधं}\ A; \skt{स्कन्ध\appbrk{}ग्र\appbrk{}ह\appbrk{}प्र\appbrk{}ति\appbrk{}षे\appbrk{}ध\appbrk{}नी\appbrk{}य\appbrk{}म\appbrk{}ध्यायं}\ H}}\skt{ }\edtext{\skt{व्याख्यास्यामः}}{\lemma{\skt{व्याख्यास्यामः}}\Afootnote{A H; \skt{व्या व्याख्यास्यामः}\ K}}\skt{~॥}\pend

\pstart\versepar \ledsidenote{\passref{6.28.3}}\skt{स्कन्द\-ग्र\-हो\-प\-सृ\-ष्टानां कुमा\-रा\-णा\-म्प्र\-श\-स्यते~। }\pend
\pstart\versepar \skt{}\edtext{\skt{वात\-घ्न\-द्रु\-म\-प\-त्राणां ः}}{\lemma{\skt{वात\appbrk{}घ्न\appbrk{}द्रु\appbrk{}म\appbrk{}प\appbrk{}त्राणां ः}}\Afootnote{K; \skt{वात\appbrk{}घ्न\appbrk{}द्रु\appbrk{}म\appbrk{}प\appbrk{}त्राणां}\ A H}}\skt{ }\edtext{\skt{क्वाथः}}{\lemma{\skt{क्वाथः}}\Afootnote{K; \skt{निष्क्वाथः}\ A; \skt{निः क्वाथः}\ H}}\skt{ परिषेचने~।}\pend

\pstart\versepar \ledsidenote{\passref{6.28.4}}\skt{तेषां मूलेषु सिद्धञ्च तैलमभ्यञ्जने हितं~। }\pend
\pstart\versepar \skt{}\edtext{\skt{अश्व\-ग\-न्ध\-सु\-रा\-मु\-स्ता\-को\-त\-र्या\-मपि}}{\lemma{\skt{अश्व\appbrk{}ग\appbrk{}न्ध\appbrk{}सु\appbrk{}रा\appbrk{}मु\appbrk{}स्ता\appbrk{}को\appbrk{}त\appbrk{}र्या\appbrk{}मपि}}\Afootnote{K H; \skt{सर्व\appbrk{}ग\appbrk{}न्ध\appbrk{}सु\appbrk{}रा\appbrk{}म\appbrk{}ण्ड\appbrk{}कै\appbrk{}ड\appbrk{}र्या\appbrk{}वाप}\ A}}\skt{ }\edtext{\skt{चेष्यते}}{\lemma{\skt{चेष्यते}}\Afootnote{K H; \skt{मिष्यते}\ A}}\skt{~।}\pend

\pstart\versepar \ledsidenote{\passref{6.28.5}}\edtext{\skt{सिद्धं}}{\lemma{\skt{सिद्धं}}\Afootnote{K H; \skt{देवदारुणि रास्नायां मधुरेषु द्रुमेषु च~। सिद्धं}\ A}}\skt{ सर्प्पिश्च सक्षीरं पान\-म\-स्यो\-प\-क\-ल्प\-येत्~। }\pend
\pstart\versepar \skt{देवदारुणि }\edtext{\skt{रास्ना च}}{\lemma{\skt{रास्ना च}}\Afootnote{K H; \emph{om.}\ A}}\skt{ }\edtext{\skt{मधु\-रे\-ष्व\-द्रि\-जेषु}}{\lemma{\skt{मधु\appbrk{}रे\appbrk{}ष्व\appbrk{}द्रि\appbrk{}जेषु}}\Afootnote{K H; \emph{om.}\ A}}\skt{ }\edtext{\skt{च}}{\lemma{\skt{च}}\Afootnote{K H; \emph{om.}\ A}}\skt{~।}\edtext{}{\lemma{\skt{पान\appbrk{}म\appbrk{}स्यो\appbrk{}प\appbrk{}क\appbrk{}ल्प\appbrk{}येत्~। देवदारुणि}}\Afootnote{K H; \skt{पानमस्मै प्रयोजयेत्}\ A}}\pend

\pstart\versepar \ledsidenote{\passref{6.28.6}}\edtext{\skt{सर्षपा}}{\lemma{\skt{सर्षपा}}\Afootnote{K H; \skt{सर्षपाः}\ A}}\skt{ }\edtext{\skt{सर्पनिर्मोकं}}{\lemma{\skt{सर्पनिर्मोकं}}\Afootnote{K H; \skt{सर्पनिर्मोको}\ A}}\skt{ }\edtext{\skt{वचां}}{\lemma{\skt{वचां}}\Afootnote{K; \skt{वचा}\ A H}}\skt{ काकादनी }\edtext{\skt{घृतं}}{\lemma{\skt{घृतं}}\Afootnote{K H; \skt{घृत}\ A}}\skt{~। }\pend
\pstart\versepar \skt{}\edtext{\edtext{\skt{उष्ट्रा\-जा\-वि\-ग\-वा\-म्वापि}}{\lemma{\skt{उष्ट्रा\appbrk{}जा\appbrk{}वि\appbrk{}ग\appbrk{}वाम्}}\Afootnote{K H; \skt{म्~। उष्ट्राजाविगवां}\ A}}}{\lemma{\skt{वापि}}\Afootnote{K H; \skt{चैव}\ A}}\skt{ रोमाण्युद्धूपनं }\edtext{\skt{भवेत्}}{\lemma{\skt{भवेत्}}\Afootnote{K H; \skt{शिशोः}\ A}}\skt{~।}\pend

\pstart\versepar \ledsidenote{\passref{6.28.7}}\skt{सोम\-व\-ल्ली\-मि\-न्द्र\-वल्लीं }\edtext{\skt{शमीविल्वाच्च}}{\lemma{\skt{शमीविल्वाच्च}}\Afootnote{K H; \skt{शमीं बिल्वस्य}\ A}}\skt{ }\edtext{\skt{वन्दकं}}{\lemma{\skt{वन्दकं}}\Afootnote{K H; \skt{कण्टकान्}\ A}}\skt{~। }\pend
\pstart\versepar \skt{मृगादन्याश्च मूलानि }\edtext{\skt{ग्रथिकान्यपि}}{\lemma{\skt{ग्रथिकान्यपि}}\Afootnote{K H; \skt{ग्रथितान्येव}\ A}}\skt{ धारयेत्~।}\pend

\pstart\versepar \ledsidenote{\passref{6.28.8}}\skt{रक्तानि माल्यानि तथा }\edtext{\skt{पताकं}}{\lemma{\skt{पताकं}}\Afootnote{K H; \skt{पताका}\ A}}\skt{ }\edtext{\skt{रक्तांश्च}}{\lemma{\skt{रक्तांश्च}}\Afootnote{K H; \skt{रक्ताश्च}\ A}}\skt{ }\edtext{\skt{गन्धां}}{\lemma{\skt{गन्धां}}\Afootnote{K H; \skt{गन्धा}\ A}}\skt{ }\edtext{\skt{विविधांश्च}}{\lemma{\skt{विविधांश्च}}\Afootnote{K H; \skt{विविधाश्च}\ A}}\skt{ }\edtext{\skt{भक्ष्यां}}{\lemma{\skt{भक्ष्यां}}\Afootnote{K H; \skt{भक्ष्याः}\ A}}\skt{~। }\pend
\pstart\versepar \skt{घण्टा च देवाय }\edtext{\skt{बलिन्निवेद्य}}{\lemma{\skt{बलिन्निवेद्य}}\Afootnote{K H; \skt{बलिर्निवेद्यः}\ A}}\skt{ }\edtext{\skt{सकुक्कुटाः}}{\lemma{\skt{सकुक्कुटाः}}\Afootnote{K H; \skt{सुकुक्कुटः}\ A}}\skt{ }\edtext{\skt{स्कन्दहते}}{\lemma{\skt{स्कन्दहते}}\Afootnote{K H; \skt{स्कन्दग्रहे}\ A}}\skt{ हिताय~।}\pend

\pstart\versepar \ledsidenote{\passref{6.28.9}}\edtext{\skt{स्नान\-स्त्रि\-रा\-त्र\-न्निशि}}{\lemma{\skt{स्नानस्त्रिरात्रन्}}\Afootnote{K H; \skt{स्नानं त्रिरात्रं}\ A}}\skt{ }\edtext{\skt{चत्वरे च}}{\lemma{\skt{चत्वरे च}}\Afootnote{K H; \skt{चत्वरेषु}\ A}}\skt{ }\edtext{\skt{कुर्या\-त्प\-रं\-शा\-लि\-य\-वै\-र्नि\-वेद्य}}{\lemma{\skt{कुर्या\appbrk{}त्प\appbrk{}रं\appbrk{}शा\appbrk{}लि\appbrk{}य\appbrk{}वै\appbrk{}र्नि\appbrk{}वेद्य}}\Afootnote{K; \skt{कुर्यात्पुरं शालि\appbrk{}य\appbrk{}वै\appbrk{}र्न\appbrk{}वैस्तु~। अद्भिश्च}\ A; \skt{कुर्या\appbrk{}त्प\appbrk{}र्शा\appbrk{}लि\appbrk{}य\appbrk{}वै\appbrk{}नि\appbrk{}वेद्य}\ H}}\skt{~। }\pend
\pstart\versepar \skt{}\edtext{\skt{गायत्रि}}{\lemma{\skt{गायत्रि}}\Afootnote{K H; \skt{गायत्र्य}\ A}}\skt{ }\edtext{\edtext{\skt{पूता\-नि\-भि\-र\-यो\-भि\-र\-ग्नि\-म्प्र\-ज्वा\-ल\-ये\-दा\-ह्र्ति\-भिश्च}}{\lemma{\skt{पूता\appbrk{}नि\appbrk{}भि\appbrk{}र\appbrk{}यो\appbrk{}भि\appbrk{}र\appbrk{}ग्निम्}}\Afootnote{K; \skt{भिमन्त्रिताभिः}\ A; \skt{पूता\appbrk{}भि\appbrk{}र\appbrk{}यो\appbrk{}भि\appbrk{}र\appbrk{}ग्निम्}\ H}}}{\lemma{\skt{प्रज्वा\appbrk{}ल\appbrk{}ये\appbrk{}दा\appbrk{}ह्र्ति\appbrk{}भिश्च}}\Afootnote{K H; \skt{प्रज्वालनं व्याहृतिभिश्च}\ A}}\skt{ }\edtext{\skt{धीमान्}}{\lemma{\skt{धीमान्}}\Afootnote{K H; \skt{वह्नेः}\ A}}\skt{~।}\pend

\pstart\versepar \ledsidenote{\passref{6.28.10}}\edtext{\skt{रक्षा\-म\-नु\-प्र\-व\-क्षामि}}{\lemma{\skt{रक्षा\appbrk{}म\appbrk{}नु\appbrk{}प्र\appbrk{}व\appbrk{}क्षामि}}\Afootnote{K H; \skt{रक्षामतः प्रवक्ष्यामि}\ A}}\skt{ }\edtext{\edtext{\skt{वाला\-ना\-म्पा\-प\-ना\-शनी}}{\lemma{\skt{वालानाम्}}\Afootnote{K H; \skt{बालानां}\ A}}}{\lemma{\skt{पापनाशनी}}\Afootnote{K H; \skt{पापनाशिनी}\ A}}\skt{~। }\pend
\pstart\versepar \skt{}\edtext{\skt{अहन्यहनि}}{\lemma{\skt{अहन्यहनि}}\Afootnote{K H; \skt{म्~। अहन्यहनि}\ A}}\skt{ कर्त्तव्या या }\edtext{\skt{भिष\-ग्नि\-र\-त\-न्द्रितैः}}{\lemma{\skt{भिष\appbrk{}ग्नि\appbrk{}र\appbrk{}त\appbrk{}न्द्रितैः}}\Afootnote{K H; \skt{भिष\appbrk{}ग्भि\appbrk{}र\appbrk{}त\appbrk{}न्द्रितैः}\ A}}\skt{~।}\pend

\pstart\versepar \ledsidenote{\passref{6.28.11}}\skt{तप\-सा\-न्ते\-ज\-सा\-ञ्चैव }\edtext{\skt{वपुषां}}{\lemma{\skt{वपुषां}}\Afootnote{K; \skt{यशसां}\ A; \skt{वपिषां}\ H}}\skt{ }\edtext{\skt{यशसान्तथा}}{\lemma{\skt{यशसान्}}\Afootnote{K H; \skt{वपुषां}\ A}}\skt{~। }\pend
\pstart\versepar \skt{}\edtext{\skt{निधानं}}{\lemma{\skt{निधानं}}\Afootnote{A K; \skt{विधानं}\ H}}\skt{ योव्ययो देवः स ते स्कन्दः प्रसीदतु~।}\pend

\pstart\versepar \ledsidenote{\passref{6.28.12}}\skt{ग्रह\-से\-ना\-प\-ति\-र्देवो देव\-से\-ना\-प\-ति\-र्विभुः~। }\pend
\pstart\versepar \skt{}\edtext{\skt{देव\-से\-ना\-रि\-पु\-हरः}}{\lemma{\skt{देव\appbrk{}से\appbrk{}ना\appbrk{}रि\appbrk{}पु\appbrk{}हरः}}\Afootnote{K H; \skt{देव\appbrk{}से\appbrk{}ना\appbrk{}ति\appbrk{}पु\appbrk{}हरः}\ A}}\skt{ पातु त्वां भगवां गुहः~।}\pend

\pstart\versepar \ledsidenote{\passref{6.28.13}}\skt{देवदेवस्य महतः पावकस्य च यः }\edtext{\skt{सुतः}}{\lemma{\skt{सुतः}}\Afootnote{A K; \skt{सुचः}\ H}}\skt{~। }\pend
\pstart\versepar \skt{गङ्गो\-मा\-कृ\-त्ति\-का\-नाञ्च स ते शर्म प्रयच्छतु~।}\pend

\pstart\versepar \ledsidenote{\passref{6.28.14}}\edtext{\skt{रक्त\-मा\-ल्या\-म्व\-र\-धरो}}{\lemma{\skt{वरधरो}}\Afootnote{K H; \skt{बरः श्रीमा}\ A}}\skt{ }\edtext{\skt{रक्त\-च\-न्द\-न\-भू\-षितः}}{\lemma{\skt{रक्त\appbrk{}च\appbrk{}न्द\appbrk{}न\appbrk{}भू\appbrk{}षितः}}\Afootnote{K H; \skt{न्रक्त\appbrk{}च\appbrk{}न्द\appbrk{}न\appbrk{}भू\appbrk{}षितः}\ A}}\skt{~। }\pend
\pstart\versepar \skt{}\edtext{\skt{रक्तदिव्यवपुः}}{\lemma{\skt{रक्तदिव्यवपुः}}\Afootnote{K H; \skt{रक्तदिव्यवपु}\ A}}\skt{ }\edtext{\skt{श्रीमां}}{\lemma{\skt{श्रीमां}}\Afootnote{K H; \skt{र्देवः}\ A}}\skt{ पातुत्वां }\edtext{\skt{क्रौञ्चसूदन}}{\lemma{\skt{क्रौञ्चसूदन}}\Afootnote{K; \skt{क्रौञ्चसूदनः}\ A H}}\skt{ }\edtext{\skt{इति}}{\lemma{\skt{इति}}\Afootnote{K H; \emph{om.}\ A}}\skt{~॥}\pend

\pstart \ledsidenote{\passref{6.28.trailer}}\edtext{\skt{कुमारे}}{\lemma{\skt{कुमारे}}\Afootnote{K; \skt{इति सुश्रु\appbrk{}त\appbrk{}सं\appbrk{}हि\appbrk{}ता\appbrk{}या\appbrk{}मु\appbrk{}त्त\appbrk{}र\appbrk{}त\appbrk{}न्त्रा\appbrk{}न्त\appbrk{}र्गते कुमारतन्त्रे स्कन्दप्रतिषेधो नामा (द्वितीयोऽध्यायः, आदि\appbrk{}तो+)\appbrk{}अ\appbrk{}ष्टा\appbrk{}विं\appbrk{}शोऽ\appbrk{}ध्यायः}\ A; \skt{स्कन्द\appbrk{}ग्र\appbrk{}ह\appbrk{}प्र\appbrk{}ति\appbrk{}षे\appbrk{}ध\appbrk{}नीयं द्वितीयोऽध्याय}\ H}}\skt{~॥ २~॥ ०~॥}\pend

\endnumbering
\beginnumbering
\eledsection*{Adhyāya 29}
\pstart \ledsidenote{\passref{6.29.1}}\skt{अथातः }\edtext{\skt{स्कन्दो\-प\-स्मा\-र\-प्र\-ति\-षे\-ध\-म\-ध्यायं}}{\lemma{\skt{स्कन्दो\appbrk{}प\appbrk{}स्मा\appbrk{}र\appbrk{}प्र\appbrk{}ति\appbrk{}षे\appbrk{}ध\appbrk{}म\appbrk{}ध्यायं}}\Afootnote{K; \skt{स्कन्दा\appbrk{}प\appbrk{}स्मा\appbrk{}र\appbrk{}प्र\appbrk{}ति\appbrk{}षेधं}\ A; \skt{स्कन्दा\appbrk{}प\appbrk{}स्मा\appbrk{}र\appbrk{}प्र\appbrk{}ति\appbrk{}षे\appbrk{}ध\appbrk{}म\appbrk{}ध्या\appbrk{}यम्}\ H}}\skt{ }\edtext{\skt{व्याख्यास्यामः}}{\lemma{\skt{व्याख्यास्यामः}}\Afootnote{A H; \skt{व्या व्याख्यास्यामः}\ K}}\skt{~॥}\pend

\pstart\versepar \ledsidenote{\passref{6.29.3}}\edtext{\skt{विल्वं}}{\lemma{\skt{विल्वं}}\Afootnote{K H; \skt{बिल्वः}\ A}}\skt{ }\edtext{\skt{शिरीषं}}{\lemma{\skt{शिरीषं}}\Afootnote{K; \skt{शिरीषो}\ A; \skt{शिरीष}\ H}}\skt{ गोलोमी }\edtext{\skt{सुर\-सा\-दि\-र्ग\-णश्च यः}}{\lemma{\skt{सुर\appbrk{}सा\appbrk{}दि\appbrk{}र्ग\appbrk{}णश्च यः}}\Afootnote{K H; \skt{सुरसादिश्च यो गणः}\ A}}\skt{~। }\pend
\pstart\versepar \skt{परिषेके प्रयोक्तव्यः }\edtext{\skt{स्कन्दा\-प\-स्मा\-र\-शा\-न्तये}}{\lemma{\skt{स्कन्दा\appbrk{}प\appbrk{}स्मा\appbrk{}र\appbrk{}शा\appbrk{}न्तये}}\Afootnote{A K; \skt{स्कन्दा\appbrk{}प\appbrk{}स्मा\appbrk{}र\appbrk{}ण\appbrk{}शा\appbrk{}न्तये}\ H}}\skt{~।}\pend

\pstart\versepar \ledsidenote{\passref{6.29.4}}\edtext{\skt{वस्तु\-मू\-त्र\-वि\-प\-क्वम्तु}}{\lemma{\skt{वस्तु\appbrk{}मू\appbrk{}त्र\appbrk{}वि\appbrk{}प\appbrk{}क्वम्तु}}\Afootnote{K H; \skt{सर्व\appbrk{}ग\appbrk{}न्ध\appbrk{}वि\appbrk{}पक्वं तु}\ A}}\skt{ तैलमभ्यञ्जने हितं~। }\pend
\pstart\versepar \skt{क्षीर\-वृ\-क्ष\-क\-षाये च काकोल्यादौ गणे तथा~।}\pend

\pstart \ledsidenote{\passref{6.29.5ab}}\skt{विपक्तव्यं घृतञ्चापि }\edtext{\skt{पाणीयम्पयसा}}{\lemma{\skt{पाणीयम्}}\Afootnote{K; \skt{पानीयं}\ A H}}\skt{ सह~।}\pend

\pstart\versepar \ledsidenote{\passref{6.29.6}}\skt{गृध्रो\-लू\-क\-पु\-री\-षाणि }\edtext{\skt{केशाहस्ते}}{\lemma{\skt{केशाहस्ते}}\Afootnote{K; \skt{केशा हस्ति}\ A H}}\skt{ }\edtext{\skt{नखं}}{\lemma{\skt{नखं}}\Afootnote{K H; \skt{नखा}\ A}}\skt{ घृतं~। }\pend
\pstart\versepar \skt{}\edtext{\skt{वृषभस्य च}}{\lemma{\skt{वृषभस्य च}}\Afootnote{K H; \skt{वृक्षभस्य च}\ A}}\skt{ रोमाणि योज्यान्युद्धूपने पि च~।}\pend

\pstart\versepar \ledsidenote{\passref{6.29.7}}\skt{अनन्तां कुक्कुटीं }\edtext{\skt{विम्वीं}}{\lemma{\skt{विम्वीं}}\Afootnote{K H; \skt{बिम्बीं}\ A}}\skt{ मर्कटीञ्चापि धारयेत्~। }\pend
\pstart\versepar \skt{पक्वापक्वानि मांसानि प्रसन्ना रुधिरं पयः~।}\pend

\pstart\versepar \ledsidenote{\passref{6.29.8}}\edtext{\skt{भूतोदनो}}{\lemma{\skt{भूतोदनो}}\Afootnote{K H; \skt{भूतौदनो}\ A}}\skt{ निवेद्यश्च }\edtext{\skt{स्कन्दापस्मारणो}}{\lemma{\skt{स्कन्दापस्मारणो}}\Afootnote{K; \skt{स्कन्दापस्मारिणेऽ}\ A; \skt{स्कन्दापस्मारिणो}\ H}}\skt{ }\edtext{\skt{वतो}}{\lemma{\skt{वतो}}\Afootnote{K; \skt{वटे}\ A; \skt{वत}\ H}}\skt{~। }\pend
\pstart\versepar \skt{}\edtext{\skt{चतुष्पथे च}}{\lemma{\skt{चतुष्पथे च}}\Afootnote{A K; \skt{चतुःपथे च}\ H}}\skt{ }\edtext{\skt{कर्तव्याः}}{\lemma{\skt{कर्तव्याः}}\Afootnote{K; \skt{कर्तव्यं}\ A H}}\skt{ स्नानमस्य यतात्मना}\pend

\pstart\versepar \ledsidenote{\passref{6.29.9}}\skt{स्कन्दा\-प\-स्मा\-र\-संज्ञो यः स्कन्दस्य दयितः सखा~। }\pend
\pstart\versepar \skt{}\edtext{\skt{विशाखसंज्ञः}}{\lemma{\skt{विशाखसंज्ञः}}\Afootnote{K H; \skt{विशाखसंज्ञ}\ A}}\skt{ }\edtext{\skt{सशिषाः}}{\lemma{\skt{सशिषाः}}\Afootnote{K; \skt{श्च शिशोः}\ A; \skt{सशिणाः}\ H}}\skt{ शिवोस्तु }\edtext{\skt{वितानेन}}{\lemma{\skt{वितानेन}}\Afootnote{K H; \skt{विकृताननः}\ A}}\skt{ }\edtext{\skt{इति}}{\lemma{\skt{इति}}\Afootnote{K H; \emph{om.}\ A}}\skt{~॥}\pend

\pstart \ledsidenote{\passref{6.29.trailer}}\edtext{\skt{कुमारे}}{\lemma{\skt{कुमारे}}\Afootnote{K; \skt{इति सुश्रु\appbrk{}त\appbrk{}सं\appbrk{}हि\appbrk{}ता\appbrk{}या\appbrk{}मु\appbrk{}त्त\appbrk{}र\appbrk{}त\appbrk{}न्त्रा\appbrk{}न्त\appbrk{}र्गते कुमारतन्त्रे स्कन्दा\appbrk{}प\appbrk{}स्मा\appbrk{}र\appbrk{}प्र\appbrk{}ति\appbrk{}षेधो नाम (तृतीयोऽध्यायः, आदितः) एको\appbrk{}न\appbrk{}त्रिं\appbrk{}शोऽ\appbrk{}ध्यायः}\ A; \skt{स्कन्दा\appbrk{}प\appbrk{}स्मा\appbrk{}र\appbrk{}प्र\appbrk{}ति\appbrk{}षे\appbrk{}ध\appbrk{}स्तृ\appbrk{}तीयः}\ H}}\skt{~॥ ३~॥ ०~॥}\pend

\endnumbering
\beginnumbering
\eledsection*{Adhyāya 30}
\pstart \ledsidenote{\passref{6.30.1}}\skt{अथातः }\edtext{\skt{शकु\-नी\-प्र\-ति\-षेधं}}{\lemma{\skt{शकु\appbrk{}नी\appbrk{}प्र\appbrk{}ति\appbrk{}षेधं}}\Afootnote{\skt{शकु\appbrk{}नी\appbrk{}प्र\appbrk{}ती\appbrk{}षेधं}\ A; \skt{शकु\appbrk{}नी\appbrk{}प्र\appbrk{}ति\appbrk{}षेधं व्या व्याख्यास्यामः}\ K; \skt{शकु\appbrk{}नि\appbrk{}प्र\appbrk{}ति\appbrk{}षे\appbrk{}ध\appbrk{}म\appbrk{}ध्यायं}\ H}}\skt{ }\edtext{\skt{व्याख्यास्यामः}}{\lemma{\skt{व्याख्यास्यामः}}\Afootnote{A H; \emph{om.}\ K}}\skt{~॥}\pend

\pstart\versepar \ledsidenote{\passref{6.30.3}}\edtext{\skt{शकु\-न्या\-भि\-प\-री\-तस्य}}{\lemma{\skt{शकु\appbrk{}न्या\appbrk{}भि\appbrk{}प\appbrk{}री\appbrk{}तस्य}}\Afootnote{K; \skt{शकु\appbrk{}न्य\appbrk{}भि\appbrk{}प\appbrk{}री\appbrk{}तस्य}\ A H}}\skt{ }\edtext{\skt{कार्यं}}{\lemma{\skt{कार्यं}}\Afootnote{K H; \skt{कार्यो}\ A}}\skt{ }\edtext{\skt{वैध्येन}}{\lemma{\skt{वैध्येन}}\Afootnote{K; \skt{वैद्येन}\ A; \skt{वैधेन}\ H}}\skt{ }\edtext{\skt{धीमता}}{\lemma{\skt{धीमता}}\Afootnote{K H; \skt{जानता}\ A}}\skt{~। }\pend
\pstart\versepar \skt{वेत\-सा\-म्र\-क\-पि\-त्थानां }\edtext{\skt{निःक्वाथः}}{\lemma{\skt{निःक्वाथः}}\Afootnote{K H; \skt{निष्क्वाथः}\ A}}\skt{ परिषेचने~।}\pend

\pstart \ledsidenote{\passref{6.30.4ab}}\edtext{\skt{कषा\-य\-म\-धु\-रै\-स्तैलैः}}{\lemma{\skt{कषा\appbrk{}य\appbrk{}म\appbrk{}धु\appbrk{}रै\appbrk{}स्तैलैः}}\Afootnote{K H; \skt{कषा\appbrk{}य\appbrk{}म\appbrk{}धु\appbrk{}रै\appbrk{}स्तैलं}\ A}}\skt{ }\edtext{\skt{कार्यञ्जाभ्यञ्जनं}}{\lemma{\skt{कार्यञ्जाभ्यञ्जनं}}\Afootnote{K H; \skt{कार्यमभ्यञ्जने}\ A}}\skt{ }\edtext{\skt{शिशाः}}{\lemma{\skt{शिशाः}}\Afootnote{K H; \skt{शिशोः}\ A}}\skt{~।}\pend

\pstart \ledsidenote{\passref{6.30.5cd}}\skt{व्रणेषूक्तानि }\edtext{\skt{चूर्पथ्यानि}}{\lemma{\skt{चूर्पथ्यानि}}\Afootnote{K; \skt{चूर्णानि पथ्यानि}\ A H}}\skt{ }\edtext{\skt{विधधीत}}{\lemma{\skt{विधधीत}}\Afootnote{K H; \skt{विविधानि}\ A}}\skt{ च~।}\pend

\pstart \ledsidenote{\passref{6.30.6ab}}\edtext{\skt{स्कन्द\-ग्र\-ह\-धू\-प\-नानि}}{\lemma{\skt{स्कन्द\appbrk{}ग्र\appbrk{}ह\appbrk{}धू\appbrk{}प\appbrk{}नानि}}\Afootnote{K H; \skt{स्कन्दग्रहे धूपनानि}\ A}}\skt{ }\edtext{\skt{पथ्यानीह}}{\lemma{\skt{पथ्यानीह}}\Afootnote{K H; \skt{तानीहा}\ A}}\skt{ }\edtext{\skt{भवन्ति च}}{\lemma{\skt{भवन्ति च}}\Afootnote{K H; \skt{पि प्रयोजयेत्}\ A}}\skt{~।}\pend

\pstart \ledsidenote{\passref{6.30.9ab}}\skt{स्कन्दा\-प\-स्मा\-र\-ण\-मनं घृतञ्चापीह पूजितं~।}\pend

\pstart \ledsidenote{\passref{6.30.7cd}}\skt{तिलतण्डुलकं }\edtext{\skt{मल्यां}}{\lemma{\skt{मल्यां}}\Afootnote{K; \skt{माल्यं}\ A H}}\skt{ }\edtext{\skt{हरि\-ता\-ल\-म्म\-नः\-शीला}}{\lemma{\skt{मनःशीला}}\Afootnote{K H; \skt{मनःशिला}\ A}}\skt{~।}\pend

\pstart\versepar \ledsidenote{\passref{6.30.8}}\skt{बलिरेष }\edtext{\skt{करञ्जैस्तु}}{\lemma{\skt{करञ्जैस्तु}}\Afootnote{K; \skt{करञ्जेषु}\ A; \skt{करञ्जे तु}\ H}}\skt{ निवेद्यो नियतात्मना~। }\pend
\pstart\versepar \skt{}\edtext{\skt{निःकुटो च}}{\lemma{\skt{निःकुटो च}}\Afootnote{K; \skt{निष्कुटे च}\ A; \skt{निष्कुटे}\ H}}\skt{ प्रयो\-क्त\-व्य\-म्स्ना\-न\-मस्य }\edtext{\skt{यथाविधिं}}{\lemma{\skt{यथाविधिं}}\Afootnote{K H; \skt{यथाविधि}\ A}}\skt{~॥}\pend

\pstart\versepar \ledsidenote{\passref{6.30.10}}\edtext{\skt{अन्तरिक्षचरा}}{\lemma{\skt{अन्तरिक्षचरा}}\Afootnote{K H; \skt{अन्तरीक्षचरा}\ A}}\skt{ देवी सर्वा\-ल\-ङ्का\-र\-भू\-षिता~। }\pend
\pstart\versepar \skt{अयोमुखी }\edtext{\skt{तीक्ष्णमुखी}}{\lemma{\skt{तीक्ष्णमुखी}}\Afootnote{K; \skt{तीक्ष्णतुण्डा शकुनी}\ A; \skt{तीक्ष्णशकुनी}\ H}}\skt{ }\edtext{\skt{देवता ते}}{\lemma{\skt{देवता ते}}\Afootnote{K H; \skt{ते}\ A}}\skt{ प्रसीदतु~।}\pend

\pstart\versepar \ledsidenote{\passref{6.30.11}}\skt{दुर्दर्शना महाकाया पिङ्गाक्षी }\edtext{\skt{भैरवस्वना}}{\lemma{\skt{भैरवस्वना}}\Afootnote{K H; \skt{भैरवस्वरा}\ A}}\skt{~। }\pend
\pstart\versepar \skt{}\edtext{\skt{लम्बो\-द\-री\-श\-ङ्कु\-कर्णी}}{\lemma{\skt{लम्बो\appbrk{}द\appbrk{}री\appbrk{}श\appbrk{}ङ्कु\appbrk{}कर्णी}}\Afootnote{A K; \skt{लंवो\appbrk{}द\appbrk{}री\appbrk{}श\appbrk{}ङ्कु\appbrk{}कर्णा}\ H}}\skt{ शकुनी ते }\edtext{\skt{प्रसीदतामिति}}{\lemma{\skt{प्रसीदतामिति}}\Afootnote{K H; \skt{प्रसीदतु}\ A}}\skt{~॥}\pend

\pstart \ledsidenote{\passref{6.30.trailer}}\edtext{\skt{कुमार~॥ ३~॥ {\symfont ❈}}}{\lemma{\skt{कुमार~॥ ३~॥ {\symfont ❈}}}\Afootnote{K; \skt{इति सुश्रु\appbrk{}त\appbrk{}सं\appbrk{}हि\appbrk{}ता\appbrk{}या\appbrk{}मु\appbrk{}त्त\appbrk{}र\appbrk{}त\appbrk{}न्त्रा\appbrk{}न्त\appbrk{}र्गते कुमारतन्त्रे शकु\appbrk{}नी\appbrk{}प्र\appbrk{}ति\appbrk{}षेधो नाम (चतुर्थोऽध्यायः, आदितः) त्रिंशोऽध्यायः}\ A; \skt{शकु\appbrk{}नी\appbrk{}प्र\appbrk{}ति\appbrk{}षे\appbrk{}ध\appbrk{}श्च\appbrk{}तुर्थः}\ H}}\skt{~॥}\pend

\endnumbering
\beginnumbering
\eledsection*{Adhyāya 31}
\pstart \ledsidenote{\passref{6.31.1}}\skt{अथातो रेव\-ती\-प्र\-ति\-षेधं }\edtext{\skt{व्याख्यास्यामः}}{\lemma{\skt{व्याख्यास्यामः}}\Afootnote{A; \emph{om.}\ K H}}\skt{~॥}\pend

\pstart\versepar \ledsidenote{\passref{6.31.3}}\edtext{\skt{अश्व\-ग\-न्धा\-ज\-शृङ्गी च}}{\lemma{\skt{अश्व\appbrk{}ग\appbrk{}न्धा\appbrk{}ज\appbrk{}शृङ्गी च}}\Afootnote{K H; \skt{अश्वगन्धा च शृङ्गी च}\ A}}\skt{ }\edtext{\skt{शारिवाथ}}{\lemma{\skt{शारिवाथ}}\Afootnote{K H; \skt{सारिवा स}\ A}}\skt{ }\edtext{\skt{पुर्नर्न्नवा}}{\lemma{\skt{पुर्नर्न्नवा}}\Afootnote{K H; \skt{पुनर्नवा}\ A}}\skt{~। }\pend
\pstart\versepar \skt{}\edtext{\skt{सहविदारी च}}{\lemma{\skt{सहविदारी च}}\Afootnote{K H; \skt{सहे तथा विदारी च}\ A}}\skt{ }\edtext{\skt{तथा}}{\lemma{\skt{तथा}}\Afootnote{K H; \emph{om.}\ A}}\skt{ }\edtext{\skt{कषायः}}{\lemma{\skt{कषायः}}\Afootnote{K H; \skt{कषायाः}\ A}}\skt{ }\edtext{\skt{परिषेचने}}{\lemma{\skt{परिषेचने}}\Afootnote{K H; \skt{सेचने हिताः}\ A}}\skt{~।}\pend

\pstart\versepar \ledsidenote{\passref{6.31.4}}\edtext{\skt{तैलमभ्यङ्गने}}{\lemma{\skt{तैलमभ्यङ्गने}}\Afootnote{K H; \skt{तैलमभ्यञ्जने}\ A}}\skt{ }\edtext{\skt{कायं}}{\lemma{\skt{कायं}}\Afootnote{K H; \skt{कार्यं}\ A}}\skt{ कुष्ठे सर्जरसे }\edtext{\skt{तथा}}{\lemma{\skt{तथा}}\Afootnote{K H; \skt{पि च}\ A}}\skt{~। }\pend
\pstart\versepar \skt{पलं }\edtext{\skt{कषायैनलदे}}{\lemma{\skt{कषायैनलदे}}\Afootnote{K H; \skt{कषायां नलदे}\ A}}\skt{ तथा गिरिकदंबके~।}\pend

\pstart\versepar \ledsidenote{\passref{6.31.5}}\edtext{\skt{वटाश्च}}{\lemma{\skt{वटाश्च}}\Afootnote{K H; \skt{धवाश्व}\ A}}\skt{ }\edtext{\skt{कर्ण्ण\-क\-कु\-भ\-धा\-त\-की\-ति\-न्दू\-केषु च}}{\lemma{\skt{कर्ण्ण\appbrk{}क\appbrk{}कु\appbrk{}भ\appbrk{}धा\appbrk{}त\appbrk{}की\appbrk{}ति\appbrk{}न्दू\appbrk{}केषु च}}\Afootnote{K H; \skt{कर्ण\appbrk{}क\appbrk{}कु\appbrk{}भ\appbrk{}धा\appbrk{}त\appbrk{}की\appbrk{}ति\appbrk{}न्दु\appbrk{}कीषु च}\ A}}\skt{~। }\pend
\pstart\versepar \skt{}\edtext{\skt{काकोल्यादौ}}{\lemma{\skt{काकोल्यादौ}}\Afootnote{K H; \skt{काकोल्यादि}\ A}}\skt{ }\edtext{\skt{गणेचापि}}{\lemma{\skt{गणेचापि}}\Afootnote{K H; \skt{गणे चैव}\ A}}\skt{ }\edtext{\skt{पाणीयं}}{\lemma{\skt{पाणीयं}}\Afootnote{K H; \skt{पानीयं}\ A}}\skt{ सर्पिरिष्यते~।}\pend

\pstart \ledsidenote{\passref{6.31.6cd}}\skt{गृध्रो\-लू\-क\-पु\-री\-षाणि }\edtext{\skt{तथा}}{\lemma{\skt{तथा}}\Afootnote{K H; \skt{यवा}\ A}}\skt{ }\edtext{\skt{यवफलाघृतं}}{\lemma{\skt{यवफलाघृतं}}\Afootnote{K H; \skt{यवफलो घृतम्}\ A}}\skt{~।}\pend

\pstart\versepar \ledsidenote{\passref{6.31.7}}\skt{संध्ययोरुभयोः कार्य\-मे\-त\-दु\-द्धू\-पनं }\edtext{\skt{शिशा}}{\lemma{\skt{शिशा}}\Afootnote{K H; \skt{शिशोः}\ A}}\skt{~। }\pend
\pstart\versepar \skt{वरु\-णा\-रि\-ष्ट\-क\-मयं }\edtext{\skt{रुमकं}}{\lemma{\skt{रुमकं}}\Afootnote{K H; \skt{रुचकं}\ A}}\skt{ }\edtext{\skt{सैन्धवं}}{\lemma{\skt{सैन्धवं}}\Afootnote{K H; \skt{सैन्दुकं}\ A}}\skt{ तथा~।}\pend

\pstart\versepar \ledsidenote{\passref{6.31.8}}\edtext{\skt{धारयेद्यापि}}{\lemma{\skt{धारयेद्यापि}}\Afootnote{K H; \skt{सततं धारयेच्चापि}\ A}}\skt{ }\edtext{\skt{सततं}}{\lemma{\skt{सततं}}\Afootnote{K H; \skt{कृतं}\ A}}\skt{ }\edtext{\skt{वचा}}{\lemma{\skt{वचा}}\Afootnote{K H; \skt{वा}\ A}}\skt{ }\edtext{\skt{पुत्रकजीवके}}{\lemma{\skt{पुत्रकजीवके}}\Afootnote{K H; \skt{पौत्रजीविकम्}\ A}}\skt{~। }\pend
\pstart\versepar \skt{शुक्लाः सुमनसो }\edtext{\skt{लाजपयः}}{\lemma{\skt{लाजपयः}}\Afootnote{K H; \skt{लाजाः पयः}\ A}}\skt{ शाल्योदनं तथा~।}\pend

\pstart\versepar \ledsidenote{\passref{6.31.9}}\edtext{\edtext{\skt{वलिन्निवेद्य}}{\lemma{\skt{वलिन्}}\Afootnote{K H; \skt{बलि}\ A}}}{\lemma{\skt{निवेद्य}}\Afootnote{K H; \skt{र्निवेद्यो}\ A}}\skt{ गोतीर्थे रेव\-त्यै\-प्र\-य\-ता\-त्मना~। }\pend
\pstart\versepar \skt{संगमे च भिषक्स्नानं कुर्या\-द्धा\-त्री\-कु\-मा\-रयोः~।}\pend

\pstart\versepar \ledsidenote{\passref{6.31.10}}\skt{नाना\-व\-स्त्र\-ध\-रा\-देवी चित्र\-मा\-ल्या\-नु\-ले\-पना~। }\pend
\pstart\versepar \skt{चलत्कुण्डलिनी श्यामा रेवतीते प्रसीदतु~।}\pend

\pstart\versepar \ledsidenote{\passref{6.31.11}}\edtext{\skt{यामु\-पा\-स\-न्नि\-स\-ततं}}{\lemma{\skt{यामु\appbrk{}पा\appbrk{}स\appbrk{}न्नि\appbrk{}स\appbrk{}ततं}}\Afootnote{K H; \skt{उपासते यां सततं}\ A}}\skt{ देव्यो विविधभूषणाः~। }\pend
\pstart\versepar \skt{}\edtext{\skt{लम्वा\-क\-रा\-ला\-वि\-नता}}{\lemma{\skt{लम्वा\appbrk{}क\appbrk{}रा\appbrk{}ला\appbrk{}वि\appbrk{}नता}}\Afootnote{K H; \skt{लम्बा कराला विनता}\ A}}\skt{ तथैव }\edtext{\skt{वहुपुत्रिका}}{\lemma{\skt{वहुपुत्रिका}}\Afootnote{K H; \skt{बहुपुत्रिका}\ A}}\skt{~। }\pend
\pstart\versepar \skt{}\edtext{\skt{रैवती}}{\lemma{\skt{रैवती}}\Afootnote{K H; \skt{रेवती}\ A}}\skt{ }\edtext{\skt{शुष्कनासा च सा च}}{\lemma{\skt{शुष्कनासा च सा च}}\Afootnote{K H; \skt{शुष्कनामा या सा ते}\ A}}\skt{ देवी }\edtext{\skt{प्रसीदतामिति}}{\lemma{\skt{प्रसीदतामिति}}\Afootnote{K H; \skt{प्रसीदतु}\ A}}\skt{~॥}\pend

\pstart \ledsidenote{\passref{6.31.trailer}}\skt{कुमार~॥ ३~॥ {\symfont ❈}~॥}\pend

\endnumbering
\beginnumbering
\eledsection*{Adhyāya 32}
\pstart \ledsidenote{\passref{6.32.1}}\skt{अथातः }\edtext{\skt{पूत\-ना\-प्र\-ति\-षेधं}}{\lemma{\skt{पूत\appbrk{}ना\appbrk{}प्र\appbrk{}ति\appbrk{}षेधं}}\Afootnote{\skt{पूत\appbrk{}ना\appbrk{}र\appbrk{}ति\appbrk{}षेधं}\ A; \skt{पूत\appbrk{}ना\appbrk{}प्र\appbrk{}ति\appbrk{}षेधं व्या व्याख्यास्यामः}\ K; \skt{पुत\appbrk{}ना\appbrk{}प्र\appbrk{}ति\appbrk{}षे\appbrk{}ध\appbrk{}म\appbrk{}ध्या\appbrk{}यम्}\ H}}\skt{ }\edtext{\skt{व्याख्यास्यामः}}{\lemma{\skt{व्याख्यास्यामः}}\Afootnote{A H; \emph{om.}\ K}}\skt{~॥}\pend

\pstart\versepar \ledsidenote{\passref{6.32.3}}\edtext{\skt{कपो\-त\-वं\-का\-ह\-रलः}}{\lemma{\skt{कपो\appbrk{}त\appbrk{}वं\appbrk{}का\appbrk{}ह\appbrk{}रलः}}\Afootnote{K; \skt{कपो\appbrk{}त\appbrk{}व\appbrk{}ङ्का\appbrk{}र\appbrk{}लुको}\ A; \skt{कपोतवंक्च अरलुर्च}\ H}}\skt{ वरुणः पारिभद्रकः~। }\pend
\pstart\versepar \skt{आस्फोता चैव योज्याः }\edtext{\skt{स्युर्बालस्य}}{\lemma{\skt{स्युर्बालस्य}}\Afootnote{K H; \skt{स्युर्बालानां}\ A}}\skt{ परिषेचने~।}\pend

\pstart\versepar \ledsidenote{\passref{6.32.4}}\skt{वचा वयः स्था गोलोमी हरितालं मनः शिला~। }\pend
\pstart\versepar \skt{कुष्ठं }\edtext{\skt{सर्जरसं}}{\lemma{\skt{सर्जरसं}}\Afootnote{K; \skt{सर्जरस}\ A H}}\skt{ }\edtext{\skt{चैव}}{\lemma{\skt{चैव}}\Afootnote{K; \skt{श्चैव}\ A H}}\skt{ तैलार्थे }\edtext{\skt{वर्गमिष्यते}}{\lemma{\skt{वर्गमिष्यते}}\Afootnote{K H; \skt{वर्ग इष्यते}\ A}}\skt{~।}\pend

\pstart\versepar \ledsidenote{\passref{6.32.5}}\skt{हितं }\edtext{\skt{घृतन्तुकाक्षीर्याः}}{\lemma{\skt{तुकाक्षीर्याः}}\Afootnote{K H; \skt{तुगाक्षीर्यां}\ A}}\skt{ सिद्धं मधुरकेषु च~। }\pend
\pstart\versepar \skt{}\edtext{\skt{कुष्थ\-ता\-ली\-स\-ख\-दिराः}}{\lemma{\skt{कुष्थ\appbrk{}ता\appbrk{}ली\appbrk{}स\appbrk{}ख\appbrk{}दिराः}}\Afootnote{K H; \skt{कुष्ठ\appbrk{}ता\appbrk{}ली\appbrk{}श\appbrk{}ख\appbrk{}दिर}\ A}}\skt{ }\edtext{\skt{र्षपज्ञने}}{\lemma{\skt{र्षपज्ञने}}\Afootnote{K; \skt{चन्दनस्यन्दने}\ A; \skt{सर्षपैष्टजने}\ H}}\skt{ तथा~।}\pend

\pstart\versepar \ledsidenote{\passref{6.32.7}}\edtext{\skt{शण\-ना\-कु\-ल\-कु\-म्भीका}}{\lemma{\skt{शण\appbrk{}ना\appbrk{}कु\appbrk{}ल\appbrk{}कु\appbrk{}म्भीका}}\Afootnote{K; \skt{गन्ध\appbrk{}ना\appbrk{}कु\appbrk{}लि\appbrk{}कु\appbrk{}म्भीके}\ A; \skt{शन\appbrk{}ना\appbrk{}कु\appbrk{}ल\appbrk{}कु\appbrk{}म्भीका}\ H}}\skt{ मज्जानो }\edtext{\skt{वदरस्य च}}{\lemma{\skt{वदरस्य च}}\Afootnote{K H; \skt{बदरस्य च}\ A}}\skt{~। }\pend
\pstart\versepar \skt{कर्कटास्थि घृतञ्चापि धूपनं }\edtext{\skt{सह}}{\lemma{\skt{सह}}\Afootnote{K H; \emph{om.}\ A}}\skt{ }\edtext{\skt{सर्षपैः}}{\lemma{\skt{सर्षपैः}}\Afootnote{K H; \skt{सर्षपैः सह}\ A}}\skt{~।}\pend

\pstart\versepar \ledsidenote{\passref{6.32.8}}\skt{काका\-द\-नी\-ञ्चि\-त्र\-फलां }\edtext{\skt{विम्वीं}}{\lemma{\skt{विम्वीं}}\Afootnote{K H; \skt{बिम्बीं}\ A}}\skt{ गुञ्जाञ्च धारयेत्~। }\pend
\pstart\versepar \skt{मत्स्योदनञ्च कुर्वीत }\edtext{\skt{कृशरम्पललं}}{\lemma{\skt{कृशरम्}}\Afootnote{K H; \skt{कृशरां}\ A}}\skt{ तथा~। }\pend
\pstart\versepar \skt{}\edtext{\edtext{\skt{शरा\-व\-स\-म्पु\-ट\-कृ\-त\-म्वलिं}}{\lemma{\skt{शरा\appbrk{}व\appbrk{}स\appbrk{}म्पु\appbrk{}ट\appbrk{}कृ\appbrk{}तम्}}\Afootnote{K H; \skt{शरावसंपुटे कृत्वा}\ A}}}{\lemma{\skt{वलिं}}\Afootnote{K H; \skt{बलिं}\ A}}\skt{ }\edtext{\skt{शून्यगृहे}}{\lemma{\skt{शून्यगृहे}}\Afootnote{A K; \skt{शून्यगृहं}\ H}}\skt{ हरेत्~।}\edtext{}{\lemma{\skt{धारयेत्~। मत्स्योदनञ्च}}\Afootnote{K H; \skt{धारयेत्~। मत्स्यौदनं च}\ A}}\pend

\pstart \ledsidenote{\passref{6.32.9ab}}\edtext{\edtext{\skt{उत्सृ\-ष्टा\-न्ना\-भि\-षे\-कश्च}}{\lemma{\skt{उत्सृष्टान्}}\Afootnote{K H; \skt{उच्छिष्टे}\ A}}}{\lemma{\skt{नाभिषेकश्च}}\Afootnote{K H; \skt{नाभिषेकेण}\ A}}\skt{ शिशोः }\edtext{\skt{स्नपनमिष्यते}}{\lemma{\skt{स्नपनमिष्यते}}\Afootnote{A K; \skt{स्नापनमिष्यते}\ H}}\skt{~॥}\pend

\pstart\versepar \ledsidenote{\passref{6.32.10}}\edtext{\skt{मलि\-ना\-म्व\-र\-स\-म्वीता}}{\lemma{\skt{मलि\appbrk{}ना\appbrk{}म्व\appbrk{}र\appbrk{}स\appbrk{}म्वीता}}\Afootnote{K H; \skt{मलि\appbrk{}ना\appbrk{}म्ब\appbrk{}र\appbrk{}सं\appbrk{}वीता}\ A}}\skt{ मलि\-ना\-रू\-क्ष\-मू\-र्धजा~। }\pend
\pstart\versepar \skt{}\edtext{\skt{शून्या\-गा\-श्र\-या\-देवी}}{\lemma{\skt{शून्या\appbrk{}गा\appbrk{}श्र\appbrk{}या\appbrk{}देवी}}\Afootnote{K; \skt{शून्यागाराश्रिता देवी}\ A; \skt{शून्या\appbrk{}गा\appbrk{}र\appbrk{}श्र\appbrk{}या\appbrk{}देवी}\ H}}\skt{ दारकं }\edtext{\skt{पातु}}{\lemma{\skt{पातु}}\Afootnote{A K; \skt{पात}\ H}}\skt{ पूतना }\edtext{\skt{इति}}{\lemma{\skt{इति}}\Afootnote{K H; \emph{om.}\ A}}\skt{~॥}\pend

\pstart \ledsidenote{\passref{6.32.trailer}}\edtext{\skt{कुमार~॥ ३~॥ {\symfont ❈}}}{\lemma{\skt{कुमार~॥ ३~॥ {\symfont ❈}}}\Afootnote{K; \skt{इति सुश्रु\appbrk{}त\appbrk{}सं\appbrk{}हि\appbrk{}ता\appbrk{}या\appbrk{}मु\appbrk{}त्त\appbrk{}र\appbrk{}त\appbrk{}न्त्रा\appbrk{}न्त\appbrk{}र्गते कुमारतन्त्रे (षष्ठोऽध्यायः, आदितः) द्वात्रिंशोऽध्यायः}\ A; \skt{पूत\appbrk{}ना\appbrk{}प्र\appbrk{}ति\appbrk{}षेधः षष्ठः}\ H}}\skt{~॥}\pend

\endnumbering
\beginnumbering
\eledsection*{Adhyāya 33}
\pstart \ledsidenote{\passref{6.33.1}}\edtext{\skt{अथा\-ता\-न्ध\-पू\-त\-ना\-प्र\-ति\-षेधं}}{\lemma{\skt{अथा\appbrk{}ता\appbrk{}न्ध\appbrk{}पू\appbrk{}त\appbrk{}ना\appbrk{}प्र\appbrk{}ति\appbrk{}षेधं}}\Afootnote{\skt{अथा\appbrk{}तोऽ\appbrk{}न्ध\appbrk{}पू\appbrk{}त\appbrk{}ना\appbrk{}प्र\appbrk{}सि\appbrk{}षेधं}\ A; \skt{अथा\appbrk{}ता\appbrk{}न्ध\appbrk{}पू\appbrk{}त\appbrk{}ना\appbrk{}प्र\appbrk{}ति\appbrk{}षेधं व्या व्याख्यास्यामः}\ K; \skt{अथा{\symfont ¦}तोऽ\appbrk{}न्ध\appbrk{}पू\appbrk{}त\appbrk{}ना\appbrk{}प्र\appbrk{}ति\appbrk{}षे\appbrk{}ध\appbrk{}म\appbrk{}ध्या\appbrk{}यम्}\ H}}\skt{ }\edtext{\skt{व्याख्यास्यामः}}{\lemma{\skt{व्याख्यास्यामः}}\Afootnote{A H; \emph{om.}\ K}}\skt{~॥}\pend

\pstart\versepar \ledsidenote{\passref{6.33.3}}\edtext{\skt{तिक्तकद्रुमाणां}}{\lemma{\skt{तिक्तकद्रुमाणां}}\Afootnote{K H; \skt{तिक्तकद्रुम}\ A}}\skt{ }\edtext{\skt{पत्रैश्च}}{\lemma{\skt{पत्रैश्च}}\Afootnote{K H; \skt{पत्राणां}\ A}}\skt{ }\edtext{\skt{क्वाथः}}{\lemma{\skt{क्वाथः}}\Afootnote{K H; \skt{कार्यः}\ A}}\skt{ }\edtext{\skt{कार्यो}}{\lemma{\skt{कार्यो}}\Afootnote{K H; \skt{क्वाथोऽ}\ A}}\skt{ वसेचने~। }\pend
\pstart\versepar \skt{सुरा सौवीरकञ्कुष्ठं हरितालम्मनः शिला~।}\pend

\pstart\versepar \ledsidenote{\passref{6.33.4}}\skt{तथा }\edtext{\skt{सर्जरसं}}{\lemma{\skt{सर्जरसं}}\Afootnote{K; \skt{सर्जरस}\ A H}}\skt{ }\edtext{\skt{चापि}}{\lemma{\skt{चापि}}\Afootnote{K; \skt{श्चैव}\ A; \skt{श्चा पि}\ H}}\skt{ तैला\-र्थ\-म\-प\-सं\-ह\-रेत्~॥ }\pend
\pstart\versepar \skt{पिप्प\-ली\-पि\-प्प\-ली\-मूलं }\edtext{\edtext{\edtext{\skt{वर्ग\-म्म\-धु\-र\-क\-म्मधुः}}{\lemma{\skt{वर्गम्}}\Afootnote{K H; \skt{वर्गो}\ A}}}{\lemma{\skt{मधुरकम्}}\Afootnote{K H; \skt{मधुरको}\ A}}}{\lemma{\skt{मधुः}}\Afootnote{K H; \skt{मधु}\ A}}\skt{~।}\edtext{}{\lemma{\skt{तैला\appbrk{}र्थ\appbrk{}म\appbrk{}प\appbrk{}सं\appbrk{}ह\appbrk{}रेत्~॥ पिप्प\appbrk{}ली\appbrk{}पि\appbrk{}प्प\appbrk{}ली\appbrk{}मूलं}}\Afootnote{K H; \skt{तैला\appbrk{}र्थ\appbrk{}मु\appbrk{}प\appbrk{}दि\appbrk{}श्यते~। पिप्पल्यः पिप्पलीमूलं}\ A}}\pend

\pstart \ledsidenote{\passref{6.33.5ab}}\edtext{\skt{सालपर्णी}}{\lemma{\skt{सालपर्णी}}\Afootnote{K H; \skt{शालपार्णी}\ A}}\skt{ }\edtext{\skt{वृहत्यौ च}}{\lemma{\skt{वृहत्यौ च}}\Afootnote{K H; \skt{बृहत्यौ च}\ A}}\skt{ }\edtext{\skt{घृतार्थमपि}}{\lemma{\skt{घृतार्थमपि}}\Afootnote{K H; \skt{घृतार्थमुप}\ A}}\skt{ }\edtext{\skt{वुद्धिमां}}{\lemma{\skt{वुद्धिमां}}\Afootnote{K; \skt{दिश्यते}\ A; \skt{वुद्धिमान्}\ H}}\skt{~।}\pend

\pstart\versepar \ledsidenote{\passref{6.33.6}}\edtext{\skt{पुरीशं}}{\lemma{\skt{पुरीशं}}\Afootnote{K H; \skt{पुरीषं}\ A}}\skt{ कौक्कुटं }\edtext{\skt{केशाः}}{\lemma{\skt{केशाः}}\Afootnote{K H; \skt{केशां}\ A}}\skt{ श्चर्म्म }\edtext{\skt{सर्व्प्पत्वगेषु च}}{\lemma{\skt{सर्व्प्पत्वगेषु च}}\Afootnote{K H; \skt{सर्पत्वचं तथा}\ A}}\skt{~। }\pend
\pstart\versepar \skt{जीर्ण्णाञ्च भिक्षुसंघाटीं }\edtext{\skt{धूप\-ना\-यो\-प\-क\-ल्प\-येत्}}{\lemma{\skt{धूप\appbrk{}ना\appbrk{}यो\appbrk{}प\appbrk{}क\appbrk{}ल्प\appbrk{}येत्}}\Afootnote{K H; \skt{धूम\appbrk{}ना\appbrk{}यो\appbrk{}प\appbrk{}क\appbrk{}ल्प\appbrk{}येत्}\ A}}\skt{~।}\pend

\pstart\versepar \ledsidenote{\passref{6.33.7}}\skt{कुक्कुटीं }\edtext{\skt{मर्क\-टी\-वि\-म्वी\-म\-न\-न्ता\-श्चापि}}{\lemma{\skt{मर्क\appbrk{}टी\appbrk{}वि\appbrk{}म्वी\appbrk{}म\appbrk{}न\appbrk{}न्ता\appbrk{}श्चापि}}\Afootnote{K; \skt{मर्कटीं शिम्बीमनन्तां चापि}\ A; \skt{मर्क्क\appbrk{}टी\appbrk{}वि\appbrk{}म्वी\appbrk{}म\appbrk{}नन्ता चापि}\ H}}\skt{ धारयेत्~॥ }\pend
\pstart\versepar \skt{मांसमामं तथा }\edtext{\skt{पक्वशोणितञ्च}}{\lemma{\skt{पक्वशोणितञ्च}}\Afootnote{K; \skt{पक्वं शोणितं च}\ A H}}\skt{ }\edtext{\skt{चतुः}}{\lemma{\skt{चतुः}}\Afootnote{K; \skt{चतु}\ A H}}\skt{ }\edtext{\skt{पथे।}}{\lemma{\skt{पथे}}\Afootnote{K; \skt{ष्पथे}\ A H}}\edtext{}{\lemma{\skt{धारयेत्~॥ मांसमामं}}\Afootnote{A K; \skt{धारयेत्~॥ मान्समासन्}\ H}}\pend

\pstart \ledsidenote{\passref{6.33.8}}\edtext{\skt{निवेद्यमन्तञ्च}}{\lemma{\skt{निवेद्यमन्तञ्च}}\Afootnote{K H; \skt{निवेद्यमन्तश्च}\ A}}\skt{ गृहे }\edtext{\skt{शिशोः}}{\lemma{\skt{शिशोः}}\Afootnote{K H; \skt{शिशो}\ A}}\skt{ }\edtext{\skt{स्नपनमिष्यते}}{\lemma{\skt{स्नपनमिष्यते}}\Afootnote{K H; \skt{रक्षानिमित्ततः~। शिशोश्च स्नपनं कुर्या\appbrk{}त्स\appbrk{}र्व\appbrk{}ग\appbrk{}न्धो\appbrk{}दकैः शुभैः}\ A}}\skt{~॥}\pend

\pstart\versepar \ledsidenote{\passref{6.33.9}}\edtext{\skt{करालपिंगला}}{\lemma{\skt{करालपिंगला}}\Afootnote{K H; \skt{कराला पिङ्गला}\ A}}\skt{ मुण्डा }\edtext{\edtext{\skt{कषा\-य\-म्व\-र\-धा\-रिणी}}{\lemma{\skt{कषायम्}}\Afootnote{K; \skt{कषायाम्}\ A; \skt{षायम्}\ H}}}{\lemma{\skt{वरधारिणी}}\Afootnote{K H; \skt{बरवासिनी}\ A}}\skt{~। }\pend
\pstart\versepar \skt{देवी बालमिमं प्रीता }\edtext{\skt{रक्ष\-ता\-म\-न्ध\-पू\-तना}}{\lemma{\skt{रक्ष\appbrk{}ता\appbrk{}म\appbrk{}न्ध\appbrk{}पू\appbrk{}तना}}\Afootnote{K H; \skt{संर\appbrk{}क्ष\appbrk{}त्व\appbrk{}न्ध\appbrk{}पू\appbrk{}तना}\ A}}\skt{ }\edtext{\skt{इति}}{\lemma{\skt{इति}}\Afootnote{K H; \emph{om.}\ A}}\skt{~॥}\pend

\pstart \ledsidenote{\passref{6.33.trailer}}\edtext{\skt{कुमार~॥ ७~॥ {\symfont ❈}}}{\lemma{\skt{कुमार~॥ ७~॥ {\symfont ❈}}}\Afootnote{K; \skt{इति सुश्रु\appbrk{}त\appbrk{}सं\appbrk{}हि\appbrk{}ता\appbrk{}या\appbrk{}मु\appbrk{}त्त\appbrk{}र\appbrk{}त\appbrk{}न्त्रा\appbrk{}न्त\appbrk{}र्गते कुमा\appbrk{}र\appbrk{}त\appbrk{}न्त्रेऽ\appbrk{}न्ध\appbrk{}पू\appbrk{}त\appbrk{}ना\appbrk{}प्र\appbrk{}ति\appbrk{}षेधो नाम (सप्तमोऽध्यायः, आदितः) त्रयस्त्रिंशोऽध्यायः}\ A; \skt{अन्ध\appbrk{}पू\appbrk{}त\appbrk{}ना\appbrk{}प्र\appbrk{}ति\appbrk{}षेधः सप्तमोध्यायः}\ H}}\skt{~॥}\pend

\endnumbering
\beginnumbering
\eledsection*{Adhyāya 34}
\pstart \ledsidenote{\passref{6.34.1}}\skt{अथातः }\edtext{\skt{शीत\-पू\-त\-ना\-प्र\-ति\-षेधं}}{\lemma{\skt{शीत\appbrk{}पू\appbrk{}त\appbrk{}ना\appbrk{}प्र\appbrk{}ति\appbrk{}षेधं}}\Afootnote{A; \skt{शीत\appbrk{}पू\appbrk{}त\appbrk{}ना\appbrk{}प्र\appbrk{}ति\appbrk{}षेधं व्या व्याख्यास्यामः}\ K; \skt{शीत\appbrk{}पू\appbrk{}त\appbrk{}ना\appbrk{}प्र\appbrk{}ति\appbrk{}षे\appbrk{}ध\appbrk{}नी\appbrk{}य\appbrk{}म\appbrk{}ध्या\appbrk{}यम्}\ H}}\skt{ }\edtext{\skt{व्याख्यास्यामः}}{\lemma{\skt{व्याख्यास्यामः}}\Afootnote{A H; \emph{om.}\ K}}\skt{~॥}\pend

\pstart\versepar \ledsidenote{\passref{6.34.3}}\edtext{\skt{विंवीकापित्थं}}{\lemma{\skt{विंवीकापित्थं}}\Afootnote{K; \skt{कपित्थं}\ A; \skt{विंवीं कपित्थ}\ H}}\skt{ }\edtext{\skt{सूरसं}}{\lemma{\skt{सूरसं}}\Afootnote{K H; \skt{सुवहां बिम्बीं}\ A}}\skt{ तथा }\edtext{\skt{विल्वप्रचीवलौ}}{\lemma{\skt{विल्वप्रचीवलौ}}\Afootnote{K H; \skt{बिल्वं प्रचीबलम्}\ A}}\skt{~। }\pend
\pstart\versepar \skt{}\edtext{\skt{नन्दी\-भ\-ल्ला\-त\-क\-ञ्चापि}}{\lemma{\skt{नन्दी\appbrk{}भ\appbrk{}ल्ला\appbrk{}त\appbrk{}कञ्}}\Afootnote{K H; \skt{नन्दीं भल्लातकं}\ A}}\skt{ परिषेके }\edtext{\skt{वचारयेत्}}{\lemma{\skt{वचारयेत्}}\Afootnote{K H; \skt{प्रयोजयेत्}\ A}}\skt{~।}\pend

\pstart\versepar \ledsidenote{\passref{6.34.4}}\edtext{\skt{वस्तमूत्रं}}{\lemma{\skt{वस्तमूत्रं}}\Afootnote{K H; \skt{बस्तमूत्रं}\ A}}\skt{ }\edtext{\skt{गवं}}{\lemma{\skt{गवं}}\Afootnote{K; \skt{गवां}\ A H}}\skt{ }\edtext{\edtext{\skt{मूत्र\-मू\-स्त\-ञ्चा\-म\-र\-दारू}}{\lemma{\skt{मूत्रमूस्तञ्}}\Afootnote{K; \skt{मूत्रं मुस्तं}\ A; \skt{मूत्रमूस्त}\ H}}\skt{ च}}{\lemma{\skt{चामरदारू च}}\Afootnote{K; \skt{च सुरदारु च}\ A; \skt{चामवदारू च}\ H}}\skt{~। }\pend
\pstart\versepar \skt{कुष्ठञ्च }\edtext{\skt{सर्वगन्धञ्च}}{\lemma{\skt{सर्वगन्धञ्च}}\Afootnote{K H; \skt{सर्वगन्धांश्च}\ A}}\skt{ तैला\-र्थ\-म\-व\-चा\-र\-येत्~।}\pend

\pstart\versepar \ledsidenote{\passref{6.34.5}}\skt{रोहि\-णी\-स\-र्ज\-ख\-दि\-र\-प\-ला\-श\-क\-कु\-भ\-त्वचः~। }\pend
\pstart\versepar \skt{}\edtext{\skt{निःक्वाथ्य}}{\lemma{\skt{निःक्वाथ्य}}\Afootnote{K H; \skt{निष्क्वाथ्य}\ A}}\skt{ }\edtext{\skt{तस्मिन्निः}}{\lemma{\skt{तस्मिन्निः}}\Afootnote{K; \skt{तस्मिन्नि}\ A H}}\skt{ }\edtext{\skt{क्वाथे}}{\lemma{\skt{क्वाथे}}\Afootnote{K H; \skt{ष्क्वाथे}\ A}}\skt{ सक्षीरं }\edtext{\skt{विपचेद्धृतं।}}{\lemma{\skt{विपचेद्धृतं}}\Afootnote{K H; \skt{विपचेद्घृतम्}\ A}}\pend

\pstart\versepar \ledsidenote{\passref{6.34.6}}\skt{गृध्रो\-लू\-क\-पु\-री\-षाणि }\edtext{\skt{वस्तगन्धां}}{\lemma{\skt{वस्तगन्धां}}\Afootnote{K H; \skt{बस्तगन्धा}\ A}}\skt{ }\edtext{\skt{सहां}}{\lemma{\skt{सहां}}\Afootnote{K H; \skt{महे}\ A}}\skt{ }\edtext{\skt{वचां}}{\lemma{\skt{वचां}}\Afootnote{K H; \skt{स्त्वचः}\ A}}\skt{~। }\pend
\pstart\versepar \skt{}\edtext{\skt{निवस्य}}{\lemma{\skt{निवस्य}}\Afootnote{K H; \skt{निम्ब}\ A}}\skt{ पत्राणि }\edtext{\skt{तथा}}{\lemma{\skt{तथा}}\Afootnote{K H; \skt{मधुकं}\ A}}\skt{ }\edtext{\skt{धूपनार्थं}}{\lemma{\skt{धूपनार्थं}}\Afootnote{A K; \skt{पूपनार्थम्}\ H}}\skt{ प्रयोजयेत्~।}\pend

\pstart\versepar \ledsidenote{\passref{6.34.7}}\skt{धारयेदपि }\edtext{\skt{गुञ्जाञ्च}}{\lemma{\skt{गुञ्जाञ्च}}\Afootnote{K H; \skt{लम्बां च}\ A}}\skt{ }\edtext{\skt{लम्वां}}{\lemma{\skt{लम्वां}}\Afootnote{K H; \skt{गुञ्जां}\ A}}\skt{ काकादनीन्तथा~। }\pend
\pstart\versepar \skt{}\edtext{\skt{नध्यां}}{\lemma{\skt{नध्यां}}\Afootnote{K; \skt{नद्यां}\ A; \skt{नध्या}\ H}}\skt{ }\edtext{\skt{मुद्गदनैश्चापि}}{\lemma{\skt{मुद्गदनैश्चापि}}\Afootnote{K; \skt{मुद्गकृतैश्चान्नै}\ A; \skt{मङ्गोदनैश्चापि}\ H}}\skt{ }\edtext{\skt{तर्प\-ये\-च्छी\-त\-पू\-तनां}}{\lemma{\skt{तर्प\appbrk{}ये\appbrk{}च्छी\appbrk{}त\appbrk{}पू\appbrk{}तनां}}\Afootnote{K H; \skt{स्तर्प\appbrk{}ये\appbrk{}च्छी\appbrk{}त\appbrk{}पू\appbrk{}त\appbrk{}नाम्}\ A}}\skt{~।}\pend

\pstart \ledsidenote{\passref{6.34.8cd}}\edtext{\skt{जलायशान्ते}}{\lemma{\skt{जलायशान्ते}}\Afootnote{K H; \skt{जलाशयान्ते}\ A}}\skt{ बालस्य }\edtext{\skt{स्पनं\-ञ्चो\-प\-दि\-श्यते}}{\lemma{\skt{स्पनंञ्}}\Afootnote{K; \skt{स्नपनं}\ A; \skt{स्पनञ्}\ H}}\skt{~।}\pend

\pstart\versepar \ledsidenote{\passref{6.34.9}}\edtext{\skt{मुग्दोदनाशिनी}}{\lemma{\skt{मुग्दोदनाशिनी}}\Afootnote{K; \skt{मुद्गौदनाशना}\ A; \skt{मुग्दौदनाशिनी}\ H}}\skt{ देवी सुरा\-शो\-णि\-त\-पा\-यिनी~। }\pend
\pstart\versepar \skt{}\edtext{\skt{जला\-श\-य\-र\-ता\-देवी}}{\lemma{\skt{जला\appbrk{}श\appbrk{}य\appbrk{}र\appbrk{}ता\appbrk{}देवी}}\Afootnote{K H; \skt{जलाशयालया देवी}\ A}}\skt{ पातु त्वां शीतपूतना }\edtext{\skt{इति}}{\lemma{\skt{इति}}\Afootnote{K H; \emph{om.}\ A}}\skt{~॥}\pend

\pstart \ledsidenote{\passref{6.34.trailer}}\edtext{\skt{कुमार}}{\lemma{\skt{कुमार}}\Afootnote{K; \skt{इति सुश्रु\appbrk{}त\appbrk{}सं\appbrk{}हि\appbrk{}ता\appbrk{}या\appbrk{}मु\appbrk{}त्त\appbrk{}र\appbrk{}त\appbrk{}न्त्रा\appbrk{}न्त\appbrk{}र्गते कुमारतन्त्रे शीत\appbrk{}पू\appbrk{}त\appbrk{}ना\appbrk{}प्र\appbrk{}ति\appbrk{}षेधो नाम (अष्टमोऽध्यायः, आदितः) चतुस्त्रिंशोऽध्यायः}\ A; \skt{शीत\appbrk{}पू\appbrk{}त\appbrk{}ना\appbrk{}प्र\appbrk{}ति\appbrk{}षे\appbrk{}धोऽ\appbrk{}ष्टमः}\ H}}\skt{~॥~॥ ०~॥}\pend

\endnumbering
\beginnumbering
\eledsection*{Adhyāya 35}
\pstart \ledsidenote{\passref{6.35.1}}\skt{अथातो }\edtext{\skt{मुण्डि\-का\-प्र\-ति\-षेधं}}{\lemma{\skt{मुण्डि\appbrk{}का\appbrk{}प्र\appbrk{}ति\appbrk{}षेधं}}\Afootnote{\skt{मुख\appbrk{}म\appbrk{}ण्डि\appbrk{}का\appbrk{}प्र\appbrk{}ति\appbrk{}षेधं}\ A; \skt{मुण्डि\appbrk{}का\appbrk{}प्र\appbrk{}ति\appbrk{}षेधं व्या व्याख्यास्यामः}\ K; \skt{मुख\appbrk{}म\appbrk{}ण्डि\appbrk{}का\appbrk{}प्र\appbrk{}ति\appbrk{}षे\appbrk{}ध\appbrk{}नी\appbrk{}य\appbrk{}म\appbrk{}ध्या\appbrk{}यम्}\ H}}\skt{ }\edtext{\skt{व्याख्यास्यामः}}{\lemma{\skt{व्याख्यास्यामः}}\Afootnote{A H; \emph{om.}\ K}}\skt{~॥}\pend

\pstart\versepar \ledsidenote{\passref{6.35.3}}\edtext{\skt{कपित्थविल्वौ}}{\lemma{\skt{कपित्थविल्वौ}}\Afootnote{K; \skt{कपित्थबिल्व}\ A; \skt{कपित्थविल्व}\ H}}\skt{ }\edtext{\skt{तर्का\-री\-वं\-श\-ग\-न्ध\-र्व\-ह\-स्तक}}{\lemma{\skt{तर्का\appbrk{}री\appbrk{}वं\appbrk{}श\appbrk{}ग\appbrk{}न्ध\appbrk{}र्व\appbrk{}ह\appbrk{}स्तक}}\Afootnote{K; \skt{तर्का\appbrk{}री\appbrk{}वां\appbrk{}शी\appbrk{}ग\appbrk{}न्ध\appbrk{}र्व\appbrk{}ह\appbrk{}स्तकाः}\ A; \skt{तर्क्का\appbrk{}री\appbrk{}वं\appbrk{}श\appbrk{}ग\appbrk{}न्ध\appbrk{}र्व्व\appbrk{}ह\appbrk{}स्तकः}\ H}}\skt{~। }\pend
\pstart\versepar \skt{}\edtext{\skt{कुवेराक्षी च}}{\lemma{\skt{कुवेराक्षी च}}\Afootnote{K H; \skt{कुबेराक्षी च}\ A}}\skt{ }\edtext{\skt{योज्यानि}}{\lemma{\skt{योज्यानि}}\Afootnote{K H; \skt{योज्याः स्यु}\ A}}\skt{ }\edtext{\skt{बालस्य}}{\lemma{\skt{बालस्य}}\Afootnote{K H; \skt{र्बालानां}\ A}}\skt{ परिषेचने~।}\pend

\pstart \ledsidenote{\passref{6.35.4}}\skt{स्वर\-सै\-र्भृ\-ङ्ग\-वृ\-क्षाणां }\edtext{\skt{तथैव}}{\lemma{\skt{तथैव}}\Afootnote{K H; \skt{तथाज}\ A}}\skt{ }\edtext{\skt{हयगन्धया}}{\lemma{\skt{हयगन्धया}}\Afootnote{K H; \skt{हरिगन्धयोः}\ A}}\skt{~। तैलं वसां च संयोज्य }\edtext{\skt{पचेदभ्यञ्जनं}}{\lemma{\skt{पचेदभ्यञ्जनं}}\Afootnote{K H; \skt{पचेदभ्यञ्जने}\ A}}\skt{ }\edtext{\skt{शिशाः}}{\lemma{\skt{शिशाः}}\Afootnote{K H; \skt{शिशोः}\ A}}\skt{~।}\pend

\pstart\versepar \ledsidenote{\passref{6.35.5}}\skt{मधूलिका }\edtext{\skt{पयस्या च}}{\lemma{\skt{पयस्या च}}\Afootnote{K H; \skt{यां पयसि}\ A}}\skt{ }\edtext{\skt{तुगाक्षीद्यो}}{\lemma{\skt{तुगाक्षीद्यो}}\Afootnote{K; \skt{तुगाक्षीर्यां}\ A; \skt{तुकाक्षीद्यो}\ H}}\skt{ }\edtext{\skt{गणौ न च}}{\lemma{\skt{गणौ न च}}\Afootnote{K; \skt{गणे तथा}\ A; \skt{गणे न च}\ H}}\skt{~। }\pend
\pstart\versepar \skt{}\edtext{\skt{मधुरैः}}{\lemma{\skt{मधुरैः}}\Afootnote{K H; \skt{मधुरे}\ A}}\skt{ }\edtext{\skt{पञ्चमूल्याश्च}}{\lemma{\skt{पञ्चमूल्याश्च}}\Afootnote{K H; \skt{पञ्चमूले च}\ A}}\skt{ }\edtext{\skt{कनीयस्या}}{\lemma{\skt{कनीयस्या}}\Afootnote{K H; \skt{कनीयसि}\ A}}\skt{ घृतम्पचेत्~।}\pend

\pstart\versepar \ledsidenote{\passref{6.35.6}}\edtext{\skt{यवाः}}{\lemma{\skt{यवाः}}\Afootnote{K H; \skt{वचा}\ A}}\skt{ सर्जरसः कुष्ठं }\edtext{\skt{सर्पि\-श्चो\-द्धू\-प\-न\-हितं}}{\lemma{\skt{सर्पि\appbrk{}श्चो\appbrk{}द्धू\appbrk{}प\appbrk{}न\appbrk{}हितं}}\Afootnote{K; \skt{सर्पिश्चोद्धूपनं हितम्}\ A H}}\skt{~। }\pend
\pstart\versepar \skt{धारयेदपि जिह्वाश्च }\edtext{\skt{चाष\-वै\-डा\-ल\-स\-त्यजाः}}{\lemma{\skt{चाष\appbrk{}वै\appbrk{}डा\appbrk{}ल\appbrk{}स\appbrk{}त्यजाः}}\Afootnote{K H; \skt{चाष\appbrk{}ची\appbrk{}र\appbrk{}ल्लि\appbrk{}स\appbrk{}र्पजाः}\ A}}\skt{~।}\pend

\pstart\versepar \ledsidenote{\passref{6.35.7}}\skt{वर्णकं चूर्णकं }\edtext{\edtext{\skt{माल्य\-म\-ञ्ज\-ना\-म्पा\-रतं}}{\lemma{\skt{माल्यमञ्जनाम्}}\Afootnote{K; \skt{माल्यमञ्जनं}\ A H}}}{\lemma{\skt{पारतं}}\Afootnote{K H; \skt{पारदं}\ A}}\skt{ तथा~। }\pend
\pstart\versepar \skt{मनः शिलां }\edtext{\skt{चोप\-ह\-रे\-त्गो\-ष्ठ\-मध्ये}}{\lemma{\skt{चोप\appbrk{}ह\appbrk{}रे\appbrk{}त्गो\appbrk{}ष्ठ\appbrk{}मध्ये}}\Afootnote{K H; \skt{चोप\appbrk{}ह\appbrk{}रे\appbrk{}द्गो\appbrk{}ष्ठ\appbrk{}मध्ये}\ A}}\skt{ }\edtext{\skt{वलिं}}{\lemma{\skt{वलिं}}\Afootnote{K H; \skt{बलिं}\ A}}\skt{ }\edtext{\skt{शुचिः}}{\lemma{\skt{शुचिः}}\Afootnote{K H; \skt{तथा}\ A}}\skt{~।}\pend

\pstart \ledsidenote{\passref{6.35.8}}\edtext{\skt{मन्त्र\-पू\-ता\-भि\-र\-द्भिश्च}}{\lemma{\skt{मन्त्र\appbrk{}पू\appbrk{}ता\appbrk{}भि\appbrk{}र\appbrk{}द्भिश्च}}\Afootnote{K H; \skt{पायसं सपुरोडाशं बल्य\appbrk{}र्थ\appbrk{}मु\appbrk{}प\appbrk{}सं\appbrk{}ह\appbrk{}रेत्~। मन्त्र\appbrk{}पू\appbrk{}ता\appbrk{}भि\appbrk{}र\appbrk{}द्भिश्च}\ A}}\skt{ तत्रैव }\edtext{\edtext{\skt{स्नपनम्भवेत्}}{\lemma{\skt{स्नपनम्}}\Afootnote{A K; \skt{स्नापनम्}\ H}}}{\lemma{\skt{भवेत्}}\Afootnote{K H; \skt{हितम्}\ A}}\skt{~॥}\pend

\pstart\versepar \ledsidenote{\passref{6.35.9}}\edtext{\skt{अलंकृता}}{\lemma{\skt{अलंकृता}}\Afootnote{A K; \skt{अलङ्कता}\ H}}\skt{ रूपवती सुभगा कामरूपिणी~। }\pend
\pstart\versepar \skt{}\edtext{\skt{गोष्ठमध्यालया}}{\lemma{\skt{गोष्ठमध्यालया}}\Afootnote{K H; \skt{गोष्ठमध्यालय}\ A}}\skt{ }\edtext{\skt{देवी}}{\lemma{\skt{देवी}}\Afootnote{K H; \skt{रता}\ A}}\skt{ पातु त्वां मुखमण्डिका }\edtext{\skt{इति}}{\lemma{\skt{इति}}\Afootnote{K H; \emph{om.}\ A}}\skt{~॥}\pend

\pstart \ledsidenote{\passref{6.35.trailer}}\edtext{\skt{कुमार}}{\lemma{\skt{कुमार}}\Afootnote{K; \skt{इति सुश्रु\appbrk{}त\appbrk{}सं\appbrk{}हि\appbrk{}ता\appbrk{}या\appbrk{}मु\appbrk{}त्त\appbrk{}र\appbrk{}त\appbrk{}न्त्रा\appbrk{}न्त\appbrk{}र्गते कुमारतन्त्रे मुख\appbrk{}म\appbrk{}ण्डि\appbrk{}का\appbrk{}प्र\appbrk{}ति\appbrk{}षेधो नाम (नवमोऽध्यायः, आदितः) पञ्च\appbrk{}त्रिं\appbrk{}श\appbrk{}त्त\appbrk{}मोऽ\appbrk{}ध्यायः}\ A; \skt{मुख\appbrk{}म\appbrk{}ण्डि\appbrk{}का\appbrk{}प्र\appbrk{}ति\appbrk{}षे\appbrk{}ध\appbrk{}नीयं नवमः}\ H}}\skt{~॥~॥ ०~॥}\pend

\endnumbering
\beginnumbering
\eledsection*{Adhyāya 36}
\pstart \ledsidenote{\passref{6.36.1}}\skt{अथातो }\edtext{\skt{नैग\-मे\-ष\-प्र\-ति\-षेधं}}{\lemma{\skt{नैग\appbrk{}मे\appbrk{}ष\appbrk{}प्र\appbrk{}ति\appbrk{}षेधं}}\Afootnote{A; \skt{नैग\appbrk{}मे\appbrk{}ष\appbrk{}प्र\appbrk{}ति\appbrk{}षेधं व्या व्याख्यास्यामः}\ K; \skt{नैग\appbrk{}मे\appbrk{}ष\appbrk{}प्र\appbrk{}ति\appbrk{}षे\appbrk{}ध\appbrk{}नी\appbrk{}य\appbrk{}म\appbrk{}ध्या\appbrk{}यम्}\ H}}\skt{ }\edtext{\skt{व्याख्यास्यामः}}{\lemma{\skt{व्याख्यास्यामः}}\Afootnote{A H; \emph{om.}\ K}}\skt{~॥}\pend

\pstart\versepar \ledsidenote{\passref{6.36.3}}\edtext{\skt{विल्वा\-ग्नि\-म\-न्थ\-पू\-तीकाः}}{\lemma{\skt{विल्वा\appbrk{}ग्नि\appbrk{}म\appbrk{}न्थ\appbrk{}पू\appbrk{}तीकाः}}\Afootnote{K H; \skt{बिल्वा\appbrk{}ग्नि\appbrk{}म\appbrk{}न्थ\appbrk{}पू\appbrk{}तीकाः}\ A}}\skt{ }\edtext{\skt{काश्मर्याः}}{\lemma{\skt{काश्मर्याः}}\Afootnote{K; \skt{कार्याः स्युः}\ A; \skt{काश्मर्या}\ H}}\skt{ परिषेचने~। }\pend
\pstart\versepar \skt{}\edtext{\edtext{\skt{सवीजञ्चापि}}{\lemma{\skt{सवीजञ्}}\Afootnote{K; \skt{सुरा सबीजं}\ A; \skt{सवीतञ्}\ H}}}{\lemma{\skt{चापि}}\Afootnote{K H; \emph{om.}\ A}}\skt{ }\edtext{\skt{धान्या\-म्ल\-म्प\-रि\-षे\-क\-प्र\-श\-स्यते}}{\lemma{\skt{परि\appbrk{}षे\appbrk{}क\appbrk{}प्र\appbrk{}श\appbrk{}स्यते}}\Afootnote{K H; \skt{परिषेके च शस्यते}\ A}}\skt{~।}\pend

\pstart\versepar \ledsidenote{\passref{6.36.4}}\skt{प्रियङ्गु सरलानन्ता शत\-पु\-ष्पा\-कु\-ट\-न्नटैः~॥ }\pend
\pstart\versepar \skt{पचेत्तैलं }\edtext{\skt{सुरा\-त\-क्र\-द\-धि\-म\-स्त्वा\-म्ल\-का\-ञ्जिकैः}}{\lemma{\skt{सुरा\appbrk{}त\appbrk{}क्र\appbrk{}द\appbrk{}धि\appbrk{}म\appbrk{}स्त्वा\appbrk{}म्ल\appbrk{}का\appbrk{}ञ्जिकैः}}\Afootnote{K; \skt{सगो\appbrk{}मू\appbrk{}त्रै\appbrk{}र्द\appbrk{}धि\appbrk{}म\appbrk{}स्त्व\appbrk{}म्ल\appbrk{}का\appbrk{}ञ्जिकैः}\ A; \skt{सुरा\appbrk{}त\appbrk{}क्र\appbrk{}द\appbrk{}धि\appbrk{}म\appbrk{}स्त्व\appbrk{}म्ल\appbrk{}का\appbrk{}ञ्जिकैः}\ H}}\skt{~।}\pend

\pstart\versepar \ledsidenote{\passref{6.36.5}}\edtext{\skt{पञ्च\-मू\-ल\-द्व\-य\-क्वाथैः}}{\lemma{\skt{पञ्च\appbrk{}मू\appbrk{}ल\appbrk{}द्व\appbrk{}य\appbrk{}क्वाथैः}}\Afootnote{K H; \skt{पञ्च\appbrk{}मू\appbrk{}ल\appbrk{}द्व\appbrk{}य\appbrk{}क्वाथे}\ A}}\skt{ क्षीरे मधुरकेषु च~। }\pend
\pstart\versepar \skt{पचेद्घृतञ्च }\edtext{\skt{मतिमां}}{\lemma{\skt{मतिमां}}\Afootnote{K H; \skt{मेधावी}\ A}}\skt{ }\edtext{\skt{खर्जूर्या}}{\lemma{\skt{खर्जूर्या}}\Afootnote{K H; \skt{खर्जूरी}\ A}}\skt{ }\edtext{\skt{मस्तकेन}}{\lemma{\skt{मस्तकेन}}\Afootnote{K H; \skt{मस्तकेऽपि}\ A}}\skt{ }\edtext{\skt{च}}{\lemma{\skt{च}}\Afootnote{K H; \skt{वा}\ A}}\skt{~।}\pend

\pstart \ledsidenote{\passref{6.36.6ab}}\edtext{\skt{वचा}}{\lemma{\skt{वचा}}\Afootnote{K H; \skt{वचां}\ A}}\skt{ }\edtext{\skt{वयस्था {\symfont ¦} {\symfont ¦} {\symfont ¦} {\symfont ¦}}}{\lemma{\skt{वयस्था {\symfont ¦} {\symfont ¦} {\symfont ¦} {\symfont ¦}}}\Afootnote{K; \skt{वयःस्थां}\ A; \skt{वयः स्था}\ H}}\skt{ }\edtext{\skt{सहितां}}{\lemma{\skt{सहितां}}\Afootnote{K; \skt{गोलोमीं}\ A; \emph{om.}\ H}}\skt{ जटिलां चापि धारयेत्~।}\pend

\pstart \ledsidenote{\passref{6.36.9}}\edtext{\skt{तिलतण्डुलकं}}{\lemma{\skt{तिलतण्डुलकं}}\Afootnote{A K; \skt{तिलतण्डुलकै}\ H}}\skt{ }\edtext{\skt{माल्यं}}{\lemma{\skt{माल्यं}}\Afootnote{A K; \skt{स्तुल्यं}\ H}}\skt{ }\edtext{\skt{भक्षांश्च}}{\lemma{\skt{भक्षांश्च}}\Afootnote{K H; \skt{भक्ष्यांश्च}\ A}}\skt{ }\edtext{\skt{विविधानि च}}{\lemma{\skt{विविधानि च}}\Afootnote{K H; \skt{विविधानपि}\ A}}\skt{~। कुमा\-र\-पि\-तृ\-मे\-षाय वृक्षमूले निवेदयेत्~।}\pend

\pstart \ledsidenote{\passref{6.36.10ab}}\edtext{\edtext{\skt{अध\-स्ता\-न्क्षी\-र\-वृ\-क्षस्य}}{\lemma{\skt{अधस्तान्}}\Afootnote{K H; \skt{अधस्ता}\ A}}}{\lemma{\skt{क्षीरवृक्षस्य}}\Afootnote{K H; \skt{द्वटवृक्षस्य}\ A}}\skt{ }\edtext{\skt{स्नप\-न\-ञ्चो\-प\-दि\-श्यते}}{\lemma{\skt{स्नपनञ्}}\Afootnote{A K; \skt{स्नापनञ्}\ H}}\skt{~।}\pend

\pstart\versepar \ledsidenote{\passref{6.36.11}}\edtext{\edtext{\skt{अजा\-न\-न\-ञ्च\-ला\-क्षिभ्रू}}{\lemma{\skt{अजाननञ्}}\Afootnote{K H; \skt{अजानन}\ A}}}{\lemma{\skt{चलाक्षिभ्रू}}\Afootnote{K H; \skt{श्चलाक्षिभ्रूः}\ A}}\skt{ कामरूपी महायशाः~। }\pend
\pstart\versepar \skt{बालं बालपिता देवो नैगमेषो }\edtext{\skt{भिरक्षतामिति}}{\lemma{\skt{भिरक्षतामिति}}\Afootnote{K H; \skt{भिरक्षतु}\ A}}\skt{~॥}\pend

\pstart \ledsidenote{\passref{6.36.trailer}}\edtext{\skt{कुमार~॥~॥ {\symfont ❈}}}{\lemma{\skt{कुमार~॥~॥ {\symfont ❈}}}\Afootnote{K; \skt{इति सुश्रु\appbrk{}त\appbrk{}सं\appbrk{}हि\appbrk{}ता\appbrk{}या\appbrk{}मु\appbrk{}त्त\appbrk{}र\appbrk{}त\appbrk{}न्त्रा\appbrk{}न्त\appbrk{}र्गते कुमारतन्त्रे नैग\appbrk{}मे\appbrk{}ष\appbrk{}प्र\appbrk{}ति\appbrk{}षेधो नाम (दशमोऽध्यायः, आदितः) षट्त्रिंशोऽध्यायः}\ A; \skt{नैग\appbrk{}मे\appbrk{}ष\appbrk{}प्र\appbrk{}ति\appbrk{}षे\appbrk{}ध\appbrk{}नीयं दशमोध्याय}\ H}}\skt{~॥}\pend

\endnumbering
\beginnumbering
\eledsection*{Adhyāya 37}
\pstart \ledsidenote{\passref{6.37.1}}\skt{अथातो }\edtext{\skt{ग्रहोत्पत्तिं}}{\lemma{\skt{ग्रहोत्पत्तिं}}\Afootnote{K; \skt{ग्रहोत्पत्ति}\ A; \skt{ग्रहोत्पत्ना}\ H}}\skt{ }\edtext{\skt{व्याख्यास्यामः}}{\lemma{\skt{व्याख्यास्यामः}}\Afootnote{K; \skt{मध्यायं व्याख्यास्यामः}\ A; \skt{माध्या\appbrk{}य\appbrk{}म्व्या\appbrk{}ख्या\appbrk{}स्यामः}\ H}}\skt{~।}\pend

\pstart \ledsidenote{\passref{6.37.2.add}}\skt{अथ खलु वत्स सुश्रुत~॥}\pend

\pstart\versepar \ledsidenote{\passref{6.37.3}}\skt{नवस्कन्दादयः प्रोक्ता बालानां य इमे ग्रहाः~। }\pend
\pstart\versepar \skt{श्रीमन्तो दिव्यवपुषो }\edtext{\skt{नारी\-पु\-रु\-ष\-वि\-ग्रहां}}{\lemma{\skt{नारी\appbrk{}पु\appbrk{}रु\appbrk{}ष\appbrk{}वि\appbrk{}ग्रहां}}\Afootnote{K; \skt{नारी\appbrk{}पु\appbrk{}रु\appbrk{}ष\appbrk{}वि\appbrk{}ग्रहाः}\ A H}}\skt{~॥}\pend

\pstart \ledsidenote{\passref{6.37.4}}\skt{एते गुहस्य र }\edtext{\skt{कृत्ति\-को\-मा\-ग्नि\-शू\-लिभिः}}{\lemma{\skt{कृत्ति\appbrk{}को\appbrk{}मा\appbrk{}ग्नि\appbrk{}शू\appbrk{}लिभिः}}\Afootnote{K; \skt{क्षार्थं कृत्ति\appbrk{}को\appbrk{}मा\appbrk{}ग्नि\appbrk{}शू\appbrk{}लिभिः}\ A H}}\skt{~। सृष्टाः }\edtext{\skt{शरवनस्थास्य}}{\lemma{\skt{शरवनस्थास्य}}\Afootnote{K H; \skt{शरवणस्थस्य}\ A}}\skt{ }\edtext{\skt{रक्षि\-त\-स\-त्म\-ते\-जसा}}{\lemma{\skt{रक्षि\appbrk{}त\appbrk{}स\appbrk{}त्म\appbrk{}ते\appbrk{}जसा}}\Afootnote{K H; \skt{रक्षि\appbrk{}त\appbrk{}स्या\appbrk{}त्म\appbrk{}ते\appbrk{}जसा}\ A}}\skt{~।}\pend

\pstart\versepar \ledsidenote{\passref{6.37.5}}\skt{स्त्रीविग्रहाग्रहा ये तु नानारूपा मयेरिताः~। }\pend
\pstart\versepar \skt{}\edtext{\skt{देवीनां}}{\lemma{\skt{देवीनां}}\Afootnote{K H; \skt{गङ्गोमा}\ A}}\skt{ कृत्तिकानान्ते भागा }\edtext{\skt{राजसतामताः}}{\lemma{\skt{राजसतामताः}}\Afootnote{K; \skt{राजसतामसाः}\ A; \skt{राजसतामसा}\ H}}\skt{~।}\pend

\pstart\versepar \ledsidenote{\passref{6.37.6}}\skt{नैगमेषस्तु }\edtext{\skt{पार्वत्याः}}{\lemma{\skt{पार्वत्याः}}\Afootnote{K; \skt{पार्वत्या}\ A H}}\skt{ सृष्टो मेषाननो ग्रहः~। }\pend
\pstart\versepar \skt{}\edtext{\skt{कुमारधारो}}{\lemma{\skt{कुमारधारो}}\Afootnote{K H; \skt{कुमारधारी}\ A}}\skt{ देवस्य }\edtext{\skt{गुह\-स्या\-त्मो\-द्भ\-वस्य}}{\lemma{\skt{गुह\appbrk{}स्या\appbrk{}त्मो\appbrk{}द्भ\appbrk{}वस्य}}\Afootnote{K H; \skt{गुहस्यात्मसमः स}\ A}}\skt{ }\edtext{\skt{वै}}{\lemma{\skt{वै}}\Afootnote{K; \skt{खा}\ A; \skt{च}\ H}}\skt{~।}\pend

\pstart\versepar \ledsidenote{\passref{6.37.7}}\edtext{\skt{स्कन्दापस्मारं}}{\lemma{\skt{स्कन्दापस्मारं}}\Afootnote{K; \skt{स्कन्दापस्मार}\ A; \skt{स्कन्दापस्मारः}\ H}}\skt{ }\edtext{\skt{सञ्जयः}}{\lemma{\skt{सञ्जयः}}\Afootnote{K H; \skt{संज्ञो यः}\ A}}\skt{ सोऽग्नि\-ना\-ग्नि\-स\-म\-द्युतिः~। }\pend
\pstart\versepar \skt{}\edtext{\skt{सैवस्कन्दसखो}}{\lemma{\skt{सैवस्कन्दसखो}}\Afootnote{K H; \skt{स च स्कन्दसखा}\ A}}\skt{ }\edtext{\skt{नाम्ना}}{\lemma{\skt{नाम्ना}}\Afootnote{K H; \skt{नाम}\ A}}\skt{ विशाख इति चोच्यते~।}\pend

\pstart\versepar \ledsidenote{\passref{6.37.8}}\skt{स्कन्दः सृष्टो भगवता देवेन त्रिपुरारिणा~। }\pend
\pstart\versepar \skt{बिभर्ति }\edtext{\skt{चापरं}}{\lemma{\skt{चापरं}}\Afootnote{K H; \skt{चापरां}\ A}}\skt{ }\edtext{\skt{संज्ञं}}{\lemma{\skt{संज्ञं}}\Afootnote{K H; \skt{संज्ञां}\ A}}\skt{ कुमार इति }\edtext{\skt{संग्रहः}}{\lemma{\skt{संग्रहः}}\Afootnote{K; \skt{स ग्रहः}\ A H}}\skt{~।}\pend

\pstart\versepar \ledsidenote{\passref{6.37.9}}\skt{बाललीलाधरो योऽयं }\edtext{\skt{देवो}}{\lemma{\skt{देवो}}\Afootnote{A K; \skt{देवं}\ H}}\skt{ }\edtext{\skt{रुद्राग्निसम्भवः}}{\lemma{\skt{रुद्राग्निसम्भवः}}\Afootnote{K H; \skt{रुद्राग्निस्ंभवः}\ A}}\skt{~। }\pend
\pstart\versepar \skt{मिथ्या }\edtext{\skt{बालेषु}}{\lemma{\skt{बालेषु}}\Afootnote{K H; \skt{चारेषु}\ A}}\skt{ }\edtext{\skt{भववां न}}{\lemma{\skt{भववां न}}\Afootnote{K H; \skt{भगवान्स्व}\ A}}\skt{ }\edtext{\skt{खल्वेषु}}{\lemma{\skt{खल्वेषु}}\Afootnote{K H; \skt{यं नैष}\ A}}\skt{ प्रवर्त्तते~।}\pend

\pstart\versepar \ledsidenote{\passref{6.37.10}}\edtext{\skt{कुमा\-र\-स्क\-न्द\-सा\-मा\-न्या\-दत्र}}{\lemma{\skt{कुमा\appbrk{}र\appbrk{}स्क\appbrk{}न्द\appbrk{}सा\appbrk{}मा\appbrk{}न्या\appbrk{}दत्र}}\Afootnote{K H; \skt{कुमारः स्कन्द\appbrk{}सा\appbrk{}मा\appbrk{}न्या\appbrk{}दत्र}\ A}}\skt{ }\edtext{\skt{केचिदपण्डिताः}}{\lemma{\skt{केचिदपण्डिताः}}\Afootnote{A; \skt{केचिद्पण्डिताः}\ K H}}\skt{~। }\pend
\pstart\versepar \skt{}\edtext{\skt{गृण्हा\-ती\-त्य\-प\-रि\-ज्ञा\-ना\-द्ब्रु\-वते}}{\lemma{\skt{गृण्हा\appbrk{}ती\appbrk{}त्य\appbrk{}प\appbrk{}रि\appbrk{}ज्ञा\appbrk{}ना\appbrk{}द्ब्रु\appbrk{}वते}}\Afootnote{K H; \skt{गृह्णा\appbrk{}ती\appbrk{}त्य\appbrk{}ल्प\appbrk{}वि\appbrk{}ज्ञाना ब्रुवते}\ A}}\skt{ देहचिन्तकाः~।}\pend

\pstart\versepar \ledsidenote{\passref{6.37.11}}\skt{ततो भगवति स्कन्दे }\edtext{\skt{देवसेनापतौ}}{\lemma{\skt{देवसेनापतौ}}\Afootnote{K H; \skt{सुरसेनापतौ}\ A}}\skt{ कृते~। }\pend
\pstart\versepar \skt{}\edtext{\skt{उप\-त\-स्थु\-र्ग्र\-हा\-स्द्युते}}{\lemma{\skt{उप\appbrk{}त\appbrk{}स्थु\appbrk{}र्ग्र\appbrk{}हा\appbrk{}स्द्युते}}\Afootnote{K H; \skt{उपतस्थुर्ग्रहाः सर्वे}\ A}}\skt{ दीप्तशक्तिधरं गुहं~॥}\pend

\pstart\versepar \ledsidenote{\passref{6.37.12}}\skt{ऊचुः }\edtext{\skt{प्राञ्जलश्चैवनं}}{\lemma{\skt{प्राञ्जलश्चैवनं}}\Afootnote{K H; \skt{प्राञ्जलयश्चैनं}\ A}}\skt{ }\edtext{\skt{वृत्तिन्नो}}{\lemma{\skt{वृत्तिन्नो}}\Afootnote{K; \skt{वृत्तिं नः}\ A; \skt{वृत्तिन्न}\ H}}\skt{ }\edtext{\skt{विध्धात्स्विति}}{\lemma{\skt{विध्धात्स्विति}}\Afootnote{\skt{संविधत्स्व वै}\ A; \skt{विध्धअत्स्विति}\ K; \skt{विदधत्स्विति}\ H}}\skt{~। }\pend
\pstart\versepar \skt{तेषामर्थे ततः स्कन्दः शिवं देव\-म\-चो\-द\-यत्~।}\pend

\pstart\versepar \ledsidenote{\passref{6.37.13}}\skt{ततो }\edtext{\skt{ग्रहां}}{\lemma{\skt{ग्रहां}}\Afootnote{A K; \skt{ग्रहा}\ H}}\skt{ }\edtext{\skt{तानुवाच}}{\lemma{\skt{तानुवाच}}\Afootnote{K; \skt{स्तानुवाच}\ A H}}\skt{ भगवां }\edtext{\skt{भगनेत्रहा}}{\lemma{\skt{भगनेत्रहा}}\Afootnote{K H; \skt{भगनेत्रहृ}\ A}}\skt{~। }\pend
\pstart\versepar \skt{}\edtext{\edtext{\skt{तैर्य\-ग्यो\-नि\-म्मा\-नु\-षञ्च}}{\lemma{\skt{तैर्यग्योनिम्}}\Afootnote{K H; \skt{त्~। तिर्यग्योनिं}\ A}}}{\lemma{\skt{मानुषञ्च}}\Afootnote{K H; \skt{मानुषं}\ A}}\skt{ }\edtext{\skt{देवञ्च}}{\lemma{\skt{देवञ्च}}\Afootnote{K H; \skt{च}\ A}}\skt{ }\edtext{\skt{त्रितयं}}{\lemma{\skt{त्रितयं}}\Afootnote{A K; \skt{तृतौअञ्}\ H}}\skt{ जगत्~।}\pend

\pstart\versepar \ledsidenote{\passref{6.37.14}}\skt{पर\-स्प\-रो\-प\-का\-रेण वर्त्तते धार्यते }\edtext{\skt{तथा}}{\lemma{\skt{तथा}}\Afootnote{K H; \skt{पि च}\ A}}\skt{~। }\pend
\pstart\versepar \skt{}\edtext{\skt{देवमनुष्यां}}{\lemma{\skt{देवमनुष्यां}}\Afootnote{K; \skt{देवा मनुष्यान्}\ A H}}\skt{ प्रीणन्ति }\edtext{\edtext{\skt{तैर्य\-ग्यो\-नि\-न्त\-थैव}}{\lemma{\skt{तैर्यग्योनिन्}}\Afootnote{K H; \skt{तैर्यग्योनीं}\ A}}}{\lemma{\skt{तथैव}}\Afootnote{K H; \skt{स्तथैव}\ A}}\skt{ च~।}\pend

\pstart\versepar \ledsidenote{\passref{6.37.15}}\edtext{\skt{यथाकालं}}{\lemma{\skt{यथाकालं}}\Afootnote{K H; \skt{वर्त\appbrk{}मा\appbrk{}नै\appbrk{}र्य\appbrk{}था\appbrk{}कालं}\ A}}\skt{ }\edtext{\skt{प्रवर्त्तन्त}}{\lemma{\skt{प्रवर्त्तन्त}}\Afootnote{K H; \skt{शीत}\ A}}\skt{ }\edtext{\skt{उष्णं}}{\lemma{\skt{उष्णं}}\Afootnote{K H; \emph{om.}\ A}}\skt{ }\edtext{\skt{वर्षहिमाद्रिभिः}}{\lemma{\skt{वर्षहिमाद्रिभिः}}\Afootnote{K H; \skt{वर्षोष्णमारुतैः}\ A}}\skt{~। }\pend
\pstart\versepar \skt{}\edtext{\skt{इज्या\-ब\-लि\-न\-म\-स्का\-रै\-र्जप्य}}{\lemma{\skt{इज्या\appbrk{}ब\appbrk{}लि\appbrk{}न\appbrk{}म\appbrk{}स्का\appbrk{}रै\appbrk{}र्जप्य}}\Afootnote{K; \skt{इज्या\appbrk{}ञ्ज\appbrk{}लि\appbrk{}न\appbrk{}म\appbrk{}स्का\appbrk{}र\appbrk{}जप}\ A; \skt{इज्यावलिनम{\symfont ¦}स्कारैर्जप्प}\ H}}\skt{ }\edtext{\skt{होमैस्तथैव}}{\lemma{\skt{होमैस्तथैव}}\Afootnote{K H; \skt{होमव्रतादि}\ A}}\skt{ }\edtext{\skt{चः}}{\lemma{\skt{चः}}\Afootnote{K; \skt{भिः}\ A; \skt{च}\ H}}\skt{~।}\pend

\pstart\versepar \ledsidenote{\passref{6.37.16}}\edtext{\skt{सम्यक्प्रयुक्तैश्च}}{\lemma{\skt{सम्यक्प्रयुक्तैश्च}}\Afootnote{K H; \skt{नराः सम्यक्प्रयुक्तैश्च}\ A}}\skt{ }\edtext{\skt{नरैः}}{\lemma{\skt{नरैः}}\Afootnote{K H; \emph{om.}\ A}}\skt{ }\edtext{\skt{प्रीणयंत्यतु}}{\lemma{\skt{प्रीणयंत्यतु}}\Afootnote{K; \skt{प्रीणन्ति त्रिदि}\ A; \skt{प्रीणयन्त्य त}\ H}}\skt{ }\edtext{\skt{देवताः}}{\lemma{\skt{देवताः}}\Afootnote{K; \skt{वेश्वरान्}\ A; \skt{देवतां}\ H}}\skt{~। }\pend
\pstart\versepar \skt{}\edtext{\skt{भागधेया}}{\lemma{\skt{भागधेया}}\Afootnote{K H; \skt{भागधेयं}\ A}}\skt{ }\edtext{\skt{विभक्ताश्च}}{\lemma{\skt{विभक्ताश्च}}\Afootnote{K H; \skt{विभक्तं च}\ A}}\skt{ शेषं किञ्चिन्न विद्यते~।}\pend

\pstart\versepar \ledsidenote{\passref{6.37.17}}\skt{तद्युष्माकं शुभा }\edtext{\skt{वृत्तिः}}{\lemma{\skt{वृत्तिः}}\Afootnote{K H; \skt{वृत्ति}\ A}}\skt{ }\edtext{\skt{पूजास्वेव}}{\lemma{\skt{पूजास्वेव}}\Afootnote{K H; \skt{र्बालेष्वेव}\ A}}\skt{ भविष्यति~। }\pend
\pstart\versepar \skt{कुलेषु येषु नेज्यंते }\edtext{\skt{देव\-ता\-पि\-त\-र\-स्तथा}}{\lemma{\skt{देव\appbrk{}ता\appbrk{}पि\appbrk{}त\appbrk{}र\appbrk{}स्तथा}}\Afootnote{K H; \skt{देवाः पितर एव च}\ A}}\skt{~।}\pend

\pstart\versepar \ledsidenote{\passref{6.37.18}}\edtext{\skt{व्राह्मणाः}}{\lemma{\skt{व्राह्मणाः}}\Afootnote{K H; \skt{ब्राह्मणाः}\ A}}\skt{ }\edtext{\skt{साधवो}}{\lemma{\skt{साधवो}}\Afootnote{K H; \skt{साधव}\ A}}\skt{ }\edtext{\skt{वापि}}{\lemma{\skt{वापि}}\Afootnote{K H; \skt{श्चैव}\ A}}\skt{ गुर\-वोऽ\-ति\-थ\-य\-स्तथा }\pend
\pstart\versepar \skt{}\edtext{\skt{निवृत्त}}{\lemma{\skt{निवृत्त}}\Afootnote{K H; \skt{निवृत्ता}\ A}}\skt{ }\edtext{\skt{शौचाचारेषु}}{\lemma{\skt{शौचाचारेषु}}\Afootnote{K H; \skt{चारशौचेषु}\ A}}\skt{ }\edtext{\skt{तथा}}{\lemma{\skt{तथा}}\Afootnote{K H; \skt{पर}\ A}}\skt{ }\edtext{\skt{क्लशोपजीविषु}}{\lemma{\skt{क्लशोपजीविषु}}\Afootnote{K H; \skt{पाकोपजीविषु}\ A}}\skt{~।}\pend

\pstart\versepar \ledsidenote{\passref{6.37.19}}\edtext{\skt{उत्स\-न्न\-भि\-क्ष\-व\-लिषु}}{\lemma{\skt{उत्स\appbrk{}न्न\appbrk{}भि\appbrk{}क्ष\appbrk{}व\appbrk{}लिषु}}\Afootnote{K H; \skt{उत्सन्नबलिभि}\ A}}\skt{ }\edtext{\skt{तथा}}{\lemma{\skt{तथा}}\Afootnote{K H; \skt{क्षेषु}\ A}}\skt{ }\edtext{\skt{भिन्नगृहेषु च}}{\lemma{\skt{भिन्नगृहेषु च}}\Afootnote{K H; \skt{भिन्न\appbrk{}कां\appbrk{}स्यो\appbrk{}प\appbrk{}भो\appbrk{}जिषु}\ A}}\skt{~। }\pend
\pstart\versepar \skt{गृहेषु }\edtext{\skt{येषु ते}}{\lemma{\skt{येषु ते}}\Afootnote{K H; \skt{तेषु ये}\ A}}\skt{ }\edtext{\skt{बालाम्स्तान्हिंसं}}{\lemma{\skt{बालाम्स्तान्हिंसं}}\Afootnote{K H; \skt{बालास्तान्गृह्णी}\ A}}\skt{ ध्वमशङ्किताः~।}\pend

\pstart\versepar \ledsidenote{\passref{6.37.20}}\skt{तत्र वो विपुला वृत्तिः पूजा चैव भविष्यति~। }\pend
\pstart\versepar \skt{एवं ग्रहाः समुत्पन्ना }\edtext{\skt{बाला}}{\lemma{\skt{बाला}}\Afootnote{K H; \skt{बालान्}\ A}}\skt{ }\edtext{\skt{हिंसन्ति}}{\lemma{\skt{हिंसन्ति}}\Afootnote{K H; \skt{गृह्णन्ति}\ A}}\skt{ चाप्यतः~।}\pend

\pstart\versepar \ledsidenote{\passref{6.37.21}}\skt{ग्रहोपसृष्टा बालास्तु }\edtext{\skt{दुश्चिकित्स्यतमाः}}{\lemma{\skt{दुश्चिकित्स्यतमाः}}\Afootnote{K H; \skt{दुश्चिकित्स्यतमा}\ A}}\skt{ }\edtext{\skt{स्मृताः}}{\lemma{\skt{स्मृताः}}\Afootnote{K H; \skt{मताः}\ A}}\skt{~। }\pend
\pstart\versepar \skt{वैक\-ल्य\-म्म\-र\-ण\-म्चापि ध्रुवं स्कन्दग्रहे }\edtext{\skt{नृणामिति}}{\lemma{\skt{नृणामिति}}\Afootnote{K H; \skt{मतम्}\ A}}\skt{~॥}\pend

\pstart \ledsidenote{\passref{6.37.21a}}\skt{कुमारभृत्ये~॥~॥ ०~॥}\pend

\pstart\versepar \ledsidenote{\passref{6.37.22}}\edtext{\skt{स्कन्दग्रहस्तु}}{\lemma{\skt{स्कन्दग्रहस्तु}}\Afootnote{K H; \skt{स्कन्दग्रहोऽत्यु}\ A}}\skt{ }\edtext{\skt{प्रथमः}}{\lemma{\skt{प्रथमः}}\Afootnote{K H; \skt{ग्रतमः}\ A}}\skt{ }\edtext{\skt{स्कन्दापस्मार}}{\lemma{\skt{स्कन्दापस्मार}}\Afootnote{K H; \skt{सर्वे}\ A}}\skt{ }\edtext{\skt{एव च}}{\lemma{\skt{एव च}}\Afootnote{K H; \skt{ष्वेव यतः}\ A}}\skt{~॥ }\pend
\pstart\versepar \skt{}\edtext{\skt{शकुनी}}{\lemma{\skt{शकुनी}}\Afootnote{K H; \skt{स्मृतः~। अन्यो वा}\ A}}\skt{ }\edtext{\skt{रेवती}}{\lemma{\skt{रेवती}}\Afootnote{K H; \skt{स}\ A}}\skt{ }\edtext{\skt{चैव}}{\lemma{\skt{चैव}}\Afootnote{K H; \skt{र्व}\ A}}\skt{ }\edtext{\skt{पूतना}}{\lemma{\skt{पूतना}}\Afootnote{K H; \skt{रूपस्तु}\ A}}\skt{ }\edtext{\skt{चान्ध}}{\lemma{\skt{चान्ध}}\Afootnote{K H; \emph{om.}\ A}}\skt{ }\edtext{\skt{पूतना}}{\lemma{\skt{पूतना}}\Afootnote{K H; \skt{न साध्यो ग्रह उच्यते}\ A}}\skt{~।}\pend

\pstart\versepar \ledsidenote{\passref{6.37.23}}\skt{पूत\-ना\-शी\-त\-नामा च तथैव मुखमण्डिका~। }\pend
\pstart\versepar \skt{नवमो नैगमेषश्च यः पित्त\-ग्र\-ह\-सं\-ज्ञितः~।}\pend

\pstart\versepar \ledsidenote{\passref{6.37.24}}\skt{प्रथ\-म\-ग्र\-ह\-लिङ्गं च अन्त च ग्रहसम्भवः~। }\pend
\pstart\versepar \skt{}\edtext{\skt{सम्पिण्डे न हि}}{\lemma{\skt{सम्पिण्डे न हि}}\Afootnote{K; \skt{सम्पिण्डु न हि}\ H}}\skt{ }\edtext{\skt{बालानामध्या}}{\lemma{\skt{बालानामध्या}}\Afootnote{K; \skt{पालानामध्या}\ H}}\skt{ }\edtext{\skt{यादमचैव}}{\lemma{\skt{यादमचैव}}\Afootnote{K; \skt{दादशचैक}\ H}}\skt{ च~॥}\pend

\pstart \ledsidenote{\passref{6.37.trailer}}\edtext{\skt{समाप्तं}}{\lemma{\skt{समाप्तं}}\Afootnote{K; \skt{इति सुश्रुतसंहिता}\ A; \emph{om.}\ H}}\skt{ }\edtext{\skt{कुमा\-र\-भृ\-त्य\-मिति}}{\lemma{\skt{कुमा\appbrk{}र\appbrk{}भृ\appbrk{}त्य\appbrk{}मिति}}\Afootnote{K; \skt{यामु\appbrk{}त्त\appbrk{}र\appbrk{}त\appbrk{}न्त्रा\appbrk{}न्त\appbrk{}र्गते कुमारतन्त्रे ग्रहोत्पत्त्यध्यायो नाम (ए\appbrk{}का\appbrk{}द\appbrk{}शोऽ\appbrk{}ध्यायः, आदितः) सप्त\appbrk{}त्रिं\appbrk{}श\appbrk{}त्त\appbrk{}मोऽ\appbrk{}ध्यायः}\ A; \skt{इति कुमारतन्त्रे ग्रहो\appbrk{}त्प\appbrk{}त्ति\appbrk{}न्ना\appbrk{}मै\appbrk{}का\appbrk{}द\appbrk{}श\appbrk{}मो\appbrk{}ध्याय उत्तरतन्त्रे पञ्च\appbrk{}त्रिं\appbrk{}श\appbrk{}त्त\appbrk{}मोऽ\appbrk{}ध्यायः}\ H}}\skt{~॥ ०~॥ ०~॥ ०~॥ ०~॥}\pend

\endnumbering
